\section{The lift of relative tropical \texorpdfstring{$2$}{2}-cycles}\label{sec:lift}

This section proves Theorem~\ref{thm:lift}. Fix an admissible $\Delta$ with a profile $\sigma$ satisfying~\ref{hyp:P}, the
configuration of Section~\ref{ssec:member}, a triangulation $\mathfrak K$ of $B=B'_\sigma$ in which $\Gamma$ is full, and a
relative tropical $2$-cycle $z\in Z(\mathfrak K)$. We construct an integral $4$-chain on the localized mirror $W^*$ of
Definition~\ref{def:Ws} whose boundary is $\sum_q\rho_q(z)\,\Sigma_q$, where $\Sigma_q\subset W^*$ is a three-sphere
isotopic to the vanishing sphere, and carry it to $Y_t$ by the transport lemma. The construction has four steps.

\begin{enumerate}
\item \emph{Transfer and normal form} (Section~\ref{ssec:lift-model}). The cycle $z$ is carried to the mirror base $B_Y$ by
the dual-pair homeomorphism of Section~\ref{sec:legendre}, and the discriminant of $B_Y$ is moved onto the positions over
which $W^*$ has exact local models (the \emph{model complex}). Near the discriminant the transferred cycle is then
replaced, without changing its boundary network, by a sum of finitely many fixed surfaces with integer coefficients: the
\emph{model germs} at the vertices of the discriminant and along its arcs, taken with the multiplicities read off from
the boundary network, and one transverse disc on each arc, with an arbitrary monodromy-invariant coefficient. What is
left lies at positive distance from the discriminant. This is the normal form of Lemma~\ref{lem:normalform}, applied on
$B_Y$, and it is the only place where $z$ enters: near the discriminant only its boundary network and the coefficients of
the transverse discs remain.
\item \emph{Pieces} (Sections~\ref{ssec:lift-legs}--\ref{ssec:lift-cross}). Each model germ and each transverse disc
has an explicit lift near the discriminant, a four-dimensional \emph{piece} of $W^*$, made of subsets of toric divisors,
of the conifold smoothing at a node, and of torus orbits through the positive real locus. At a node the pieces of the
two passages contain the same sphere $\Sigma_q$, and Lemma~\ref{lem:SI} computes their boundary there.
\item \emph{The regular part} (Sections~\ref{ssec:lift-collar} and~\ref{ssec:lift-regular}). The rest of the normalised
chain is lifted over an abstract $2$-complex $\mathscr S$: the disjoint union of its \emph{summands} (the germs, the
discs, the remainder), glued only along the curves where their boundaries cancel. Over $\mathscr S$ the lift is a
bundle of two-tori with integer multiplicities, completed by $3$-chains in the fibres over the edges where several
summands meet. Summands that meet transversally in $B_Y$ are not identified in $\mathscr S$; their lifts are
separate chains, which add in $W^*$. The map to $W^*$ is built cell by cell in regions of $W^*$ that are aspherical,
with fundamental group the local lattice of the base.
\item \emph{The fibre class} (Proposition~\ref{lem:F}). The boundary of the resulting chain is $\sum_q\rho_q\Sigma_q$ plus
closed $3$-cycles in fibres over finitely many points. Each is a multiple of the class of a fibre three-torus, and
pairing with the positive real locus of $W^*$, which meets a fibre torus once and misses every $\Sigma_q$, shows that
the multiples sum to zero.
\end{enumerate}

So the lift is built summand by summand on an abstract domain; summands are glued only along curves where
their boundaries cancel; its boundary, $\sum_q\rho_q(z)\Sigma_q$, is additive in $z$; and the only global input is
Proposition~\ref{lem:F}. The strict tropical
$2$-cycles of \cite{CBM} are a special case, in which every multiplicity of a model germ is $0$ or $\pm1$
(Section~\ref{ssec:lift-cbm}).

The local terms used below refer to the following geometric objects. On $B_Y$, the images of the negative
vertices of $B'_\sigma$ have positive type; these are the \emph{model vertices}. The images of the positive
vertices are the \emph{turning points} (Proposition~\ref{prop:dict}).
\begin{center}\small
\begin{tabular}{p{3.4cm}p{10.8cm}}
term & meaning in the construction\\\hline
model complex & the mirror base with the discriminant moved to the positions of the exact local equations\\
model germ & one of the fixed local surfaces used in the normal form near a discriminant vertex or arc\\
summand & a germ, transverse disc or part of the remainder, with its integral coefficient\\
depth & the nonnegative transverse coordinate measuring distance from the discriminant within a summand\\
inner boundary & the frontier where a local piece joins the lift over the regular part\\
collar & the family of paths, together with its torus orbits, joining these two parts\\
region & an open subset of the localized mirror homotopy equivalent to the torus of the local lattice\\
fibre complex & the cell complex of tori and filling chains used over a simplex of the abstract domain
\end{tabular}
\end{center}

Throughout the section $c$ is a real coefficient vector with $c_0>0>c_n$ for $n\ne0$, in the window of
Definition~\ref{def:window}, with $r\le r_1$ and $\eta_0\le r^2/20$, and $t<t_0(r,\eta_0,\mathscr C_0)$; these are the
members of Theorem~\ref{thm:main} (Section~\ref{ssec:member}). We use the notation of Section~\ref{sec:regime}: $L=\log(1/t)$, the valuations
$\operatorname{val}_n$, the deficits, the weights $\psi_n$, the cut-off constants $\theta_0,\dots,\theta_3$ and $D$,
the constants $e_{\max}$ and $c_M$, the node data of Definition~\ref{def:node}, and the identification of $B_Y$ with the
union of the bounded cells of the tropical hypersurface $\operatorname{Trop}\subset M_\R$. The lattices $\Lambda_Y$,
$\Lambda^M$ are lattices of integral tangent vectors of $B_Y$, contained in $M$; relative cycles on $B_Y$ have
coefficients in the dual local system, the covectors, whose values on $\operatorname{int}F^*_n$ are the classes in
$N/\Z n$ of the elements of $N$.

\subsection{The model complex, the transferred cycle and its normal form}\label{ssec:lift-model}

We first move the discriminant to the positions of the exact local models and transfer the cycle to the mirror base.
The normal-form lemma then separates its boundary germs, transverse discs and a remainder disjoint from the discriminant.

\subsubsection*{The faces \texorpdfstring{$E_e$}{Ee} and the model positions}
The local models of $W^*$ sit over definite points of $B_Y$: the Morse chart over the midpoint of the edge $J_P$ of
Lemma~\ref{lem:gaps}, the models at the vertices of positive type over the segments $J^\flat_\omega$ of
Lemma~\ref{lem:windows}(3), and the leg models over the interior polygons $\Omega_e$ of Definition~\ref{def:legcell}. In general the discriminant $\Gamma_Y$
does not pass through these positions (Remark~\ref{rem:literalpos}). We move it there by a homeomorphism of $B_Y$ that
preserves the dual-pair complex of Proposition~\ref{prop:L1}; Section~\ref{sec:legendre} uses only the combinatorics of
the cells of that complex, not their positions inside the faces of $B_Y$.

We use the faces $E_e$, the interior polygons $\Omega_e$ and the turning points of Definition~\ref{def:legcell} and the
paragraph after it.

For every edge $[u,u']$ of $\mathcal T$ in the dual edge $G^\circ$ of a two-face $G$, or in a two-face $E^\vee$ of
$\Delta^\circ$ dual to an edge $E$ of $\Delta$, we take the vector $w_u\in N$ of Section~\ref{ssec:member}, with
$\langle u,w_u\rangle=1$ and $\langle u',w_u\rangle=0$, put $w_{u'}:=-n_0-w_u$ for a lattice point $n_0$ of the face
$[u,u']^\vee$ of $\Delta$, and write $\alpha_{uu'}:=\langle\,\cdot\,,w_u-w_{u'}\rangle$; at a node these are the vectors
of~\eqref{eq:nodeexp}, with $n_0=a$. At the tropical position of a point of the torus,
$L\alpha_{uu'}=\log|Z|-\log|Z'|$ for the characters $Z=\chi^{w_u}$, $Z'=\chi^{w_{u'}}$, which are the Cox coordinates of
$u$ and $u'$ times units on the chart $U_{u,u'}$. Since $\langle u'-u,w_u-w_{u'}\rangle=-2$, on every polygon $E_e$ with
$e\subseteq[u,u']^\vee$ the level lines of $\alpha_{uu'}$ are parallel lines transverse to $u'-u$, and the edges
$J_\omega$ of $E_e$ for the two-cells $\omega\supset e$ with $\omega^\vee\ni u,u'$ are parallel to $u'-u$.

\begin{definition}[model positions]\label{def:modelpos}
\begin{enumerate}
\item The \emph{model node} of a node parallelogram $P$ is the midpoint $m_P$ of $J_P$. For a unimodular triangle $\omega$
of $\mathcal F_C$ in a two-face $G$, let $u_0,\dots,u_k$ be the lattice points of $G^\circ$ in order, and let $J^\flat_\omega\subseteq J_\omega$
be the segment of Lemma~\ref{lem:windows}(3), which contains the points $p_\omega+s(u_k-u_0)$ with $|s-s_{\rm mid}|\le
1/(4e_{\max})$. The \emph{model vertices} of $\omega$ are points
$m_{\omega,1},\dots,m_{\omega,k}$ of $J^\flat_\omega$ at mutual distances at least $4\theta_3$, in the order in which the
zero-cells $C_Y([u_{i-1},u_i],\omega)$ lie on $J_\omega$. They exist, since that segment has length at least
$k\,c_M/(2e_{\max})\ge4(k-1)\theta_3$.
\item Let $e$ be an edge of $\mathcal F_C$ in the two-skeleton with $\dim e^\vee\ge1$, $\omega\supset e$ a two-cell of a
two-face $G$, and $m$ a model node or vertex of $\omega$, with edge $[u,u']$ of $\mathcal T$. A \emph{leaf path} from $m$
is a path in $E_e$ made of two segments $[m,m']$ and $[m',p_e]$, where $[m,m']$ lies on the level line of $\alpha_{uu'}$
through $m$ and has positive length at most $\theta_3$, and $p_e\in\Omega_e$. Since $\theta_3\le\theta_2/4$, it is a path
of Lemma~\ref{lem:windows}(1$'$). Let $\theta_4\in(0,\theta_2/2]$ be a constant, depending only on the
configuration and the leaf paths, such that at every model node and vertex the points of the leaf paths whose torus
coordinates are within $2\theta_4$ of those of $m$ lie on the first segments.
\item A \emph{model arc system} of $e$ is a family of piecewise linear arcs in $E_e$, each meeting every line parallel to
the direction $u'-u$ of its edge $[u,u']$ of $\mathcal T$ at most once, as follows. If $e$ lies in the interior of a
two-face $G$ with $k$ edges of $\mathcal T$ in $G^\circ$, it consists of $k$ disjoint arcs; the $i$-th joins the $i$-th
model points on the two edges $J_{\omega_\pm}$ of $E_e$, where $\omega_\pm$ are the two two-cells of $G$ containing $e$,
and consists of a leaf path from each of them and an arc in $\Omega_e$. If $e$ lies on an edge $E$ of $\Delta$, it is a
plane graph properly embedded in $E_e$, isomorphic, as a plane graph with its leaves on $\partial E_e$ in their cyclic
order, to $\Gamma_Y\cap E_e$, the graph dual to the triangulation of $E^\vee$ by $\mathcal T$ together with one leaf for
each edge of $\mathcal T$ on $\partial E^\vee$ (Proposition~\ref{prop:L2}(2)); its leaves are the model points on the
edges $J_\omega$ of $E_e$, its leaf edges begin with leaf paths, and its interior vertices and the remaining parts of
its edges lie in $\Omega_e$, each edge dual to $[u_i,u_j]$ being transverse to $u_j-u_i$.
\end{enumerate}
\end{definition}

In~(3) the graph is a tree exactly when $E^\vee$ has no lattice point in its interior; this fails on each of the
examples $X_9$, $X_{19}$ and $X_{20}$ of Section~\ref{sec:examples}. When $e$ lies in the interior of a two-face and both two-cells
$\omega_\pm$ are node parallelograms, the arc joins two model nodes. This occurs at a pentagon (the edge from the centre
shared by its two parallelograms) and at a hexagon on the $A_2$ branch, where the three rhombi share pairwise the edges
$[o,w_{j_0}]$, $[o,w_{j_0+2}]$, $[o,w_{j_0+4}]$ and the three arcs between their model nodes form a cycle of $\Gamma^M$
without vertices of valency three. Nothing below distinguishes these arcs from others.

\subsubsection*{Model arc systems}
We construct model arc systems as homothetic copies of $\Gamma_Y\cap E_e$, joined to the model points by leaf paths. The
arcs, vertices and leaves of $\Gamma_Y\cap E_e$ itself, in their actual positions, are called \emph{original}. Fix $e$ as
in Definition~\ref{def:modelpos}(3). In the picture of $\operatorname{Trop}$ the face $E_e$ is
$-(e^\vee+\nabla^{H'}_e)$, where $e^\vee+\nabla^{H'}_e$ is the face $F^*_e$ of $\partial\nabla^H$ and $\nabla^{H'}_e$ is a
polygon parallel to $e^\vee$. By Lemma~\ref{lem:cayley}(2) for $B_Y$, the cells of $\mathscr P_Y$ in $E_e$ are the sets
$-(F_1(\zeta)+F_2(\zeta))$, $\zeta\in N_\R$, where $F_1(\zeta)$ is the cell of $\mathcal T$ in $e^\vee$ on which
$a\mapsto\langle a,\zeta\rangle+\check h(a)$ is minimal and $F_2(\zeta)$ is the face of $\nabla^{H'}_e$ on which $\zeta$
is minimal; $F_1(\zeta)$ is the lower support of the cell, and only the class of $\zeta$ modulo $\R n+\R n'$ matters. For
a cell $\mathfrak c$ we write $b_{\mathfrak c}$ for the average of its vertices, the vertex of the barycentric subdivision.

\begin{lemma}[original arcs]\label{lem:litarcs}
Let $\tau=[u_i,u_j]$ be an edge of $\mathcal T$ in $e^\vee$, and let $y\in N_\R$ satisfy $\langle u_j-u_i,y\rangle=0$ and
$y\notin\R n+\R n'$. Then the closure of the arc $C_Y(\tau,e)$ is a polygonal arc along which $\langle\,\cdot\,,y\rangle$
is strictly monotone. In particular it meets every line parallel to $u_j-u_i$ at most once.
\end{lemma}

The proof is in Appendix~\ref{app:liftproofs}.

\emph{The construction.} Fix $o_e\in\Omega_e$ and $\mathfrak t_e\in(0,1/2]$ with $\mathfrak k_e(E_e)\subset\Omega_e$ for the
homothety $\mathfrak k_e(p):=o_e+\mathfrak t_e(p-o_e)$. For an edge $J$ of $E_e$ let $y_J$ be the affine function on the
plane of $E_e$ with $y_J=0$ on $J$ and $y_J(o_e)=1$, so that $y_J\ge0$ on $E_e$ and $y_J\circ\mathfrak
k_e=1-\mathfrak t_e+\mathfrak t_ey_J$. Let $S_J:=\{o_e+t(b-o_e):b\in\operatorname{ri}J,\ t>0\}$, an open convex sector; the sectors
of distinct edges are disjoint. A model point $m$ on $J=J_\omega$, with edge $[u,u']$ of $\mathcal T$, lies in
$\operatorname{ri}J\subset S_J$, and the level line of $\alpha_{uu'}$ through $m$ is transverse to $J$. Choose
$\eta_J\in(0,1/4]$ so small that for every model point $m$ on $J$ the point $m'$ of that level line with $y_J(m')=\eta_J$
satisfies $|m'-m|\le\theta_3$ and $m'\in S_J$. The \emph{leaf path} of $m$ in $E_e$ is $[m,m']\cup[m',\mathfrak
k_e(\bar m)]$, where $\bar m$ is the original leaf $C_Y([u,u'],\omega)$ corresponding to $m$ (the $i$-th model point to
the $i$-th original leaf). It is a leaf path in the sense of Definition~\ref{def:modelpos}(2), since
$\mathfrak k_e(\bar m)\in\Omega_e$.

\begin{lemma}[leaf paths]\label{lem:leafpaths}
The segments of the leaf paths in $E_e$ are pairwise disjoint, except that the two segments of one path share the
point $m'$. Each leaf path meets $\partial E_e$ only at its model point, and $\mathfrak k_e(E_e)$ only at its end
$\mathfrak k_e(\bar m)$.
\end{lemma}

The proof is in Appendix~\ref{app:liftproofs}.

\begin{lemma}[model arc systems]\label{lem:arcsystems}
For every triangle ${s_\triangle}$ of $\mathcal T$ in $e^\vee$ let $\mathfrak c_{s_\triangle}$ be the polygon of the proof of
Lemma~\ref{lem:litarcs}, let $\mathfrak c$ run over the three edges of $\mathfrak c_{s_\triangle}$ at which the closed arcs
$\overline{C_Y(\tau,e)}$, $\tau\subset{s_\triangle}$, end with the segments $[b_{\mathfrak c},b_{\mathfrak c_{s_\triangle}}]$, and let
$v_{s_\triangle}$ be the centroid of the triangle with vertices the three points $b_{\mathfrak c}$. Let $\Gamma'_e$ be
$\Gamma_Y\cap E_e$ with each segment $[b_{\mathfrak c},b_{\mathfrak c_{s_\triangle}}]$ replaced by $[b_{\mathfrak c},v_{s_\triangle}]$,
and $\mathscr G_e$ the union of $\mathfrak k_e(\Gamma'_e)$ and of the leaf paths in $E_e$. Then $\mathscr G_e$ is a model
arc system of $e$, and the correspondence which maps $C_Y({s_\triangle},e)$ to $\mathfrak k_e(v_{s_\triangle})$ and each original leaf
to its model point is an isomorphism of plane graphs from $\Gamma_Y\cap E_e$ onto $\mathscr G_e$, preserving the
rotation at every vertex and the cyclic order of the leaves on $\partial E_e$ (the points $\mathfrak k_e(\bar m)$ being
points of degree two). At every interior vertex $\mathfrak k_e(v_{s_\triangle})$ the three angles between consecutive edges
are less than $\pi$.
\end{lemma}

The proof is in Appendix~\ref{app:liftproofs}.

\subsubsection*{The model complex}

\begin{lemma}[the model complex]\label{lem:modelcx}
Fix model positions and model arc systems. There are a regular cell decomposition $\{C^M({\mathfrak p})\}_{{\mathfrak p}\in\mathrm{DP}}$ of $B_Y$
with face poset $(\mathrm{DP},\preceq)$ and a piecewise linear homeomorphism $\Upsilon^M$ of $B_Y$ such that:
\begin{enumerate}
\item $\Upsilon^M(C_Y({\mathfrak p}))=C^M({\mathfrak p})$ for every dual pair ${\mathfrak p}$; $\Upsilon^M$ maps every face $F^*_\omega$ of $B_Y$ onto
itself, and it is isotopic to the identity through homeomorphisms with this property;
\item the zero-cells $C^M([u,u'],\omega)$, $\omega$ a two-cell of a two-face, are the model nodes and model vertices, the
zero-cells $C^M(\tau,e)$, $\tau$ a triangle of $\mathcal T$, are the interior vertices of the model arc systems, and the
arcs $C^M([u,u'],e)$ are the arcs and edges of the model arc systems;
\item the local system $\Lambda^M:=\Upsilon^M_*\Lambda_Y$ is the constant lattice $n^\perp\cap M$ on $\operatorname{int}F^*_n$
and $M/\Z u$ on $R^M_u:=\bigcup_\omega C^M(u,\omega)$, glued on $R^M_u\cap\operatorname{int}F^*_n=C^M(u,n)$ by the
projection $n^\perp\cap M\to M/\Z u$, as $\Lambda_Y$ is (Proposition~\ref{prop:L3}).
\end{enumerate}
\end{lemma}

The proof is in Appendix~\ref{app:liftproofs}.

We write $\Gamma^M:=\Upsilon^M(\Gamma_Y)$ for the \emph{model discriminant}; its vertices are the model nodes, the model
vertices and the turning points $C^M(\tau,e)$, and we call them the \emph{$\Gamma^M$-vertices}. The model vertices are of
positive type and the turning points of negative type on $B_Y$ (Proposition~\ref{prop:L2}(3)).

\begin{remark}[the original positions]\label{rem:literalpos}
The node $C_Y([u,u'],P)$ of $\Gamma_Y$ is the barycentre of the cell of $\mathscr P_Y$ on $J_P$ with lower support
$[u,u']$. If $\check h(u)\ne\check h(u')$ this cell is a segment of length one at an end of $J_P$, and the node lies at
distance $(\mathfrak g_P-1)/2$ from $m_P$, where the Morse chart of $W^*$ sits. This is why the model complex is
needed.
\end{remark}

\subsubsection*{Charts at the model discriminant}
At a model node $m_P$ we use the tropical chart $T_q=(y_1,y_2,\alpha_q)$ of Lemma~\ref{lem:nodemodel}(4), with
$\alpha_q=\alpha_{uu'}-\alpha_{uu'}(m_P)$. Along the model arcs we use the following charts.

\begin{lemma}[tropical charts along the model arcs]\label{lem:tropcharts}
\begin{enumerate}
\item Let $\omega$ be a two-cell of a two-face $G$, $[u,u']$ an edge of $\mathcal T$ in $G^\circ$ and $m$ its model point on
$J_\omega$. Near $m$, $C^M(u,\omega)$ and $C^M(u',\omega)$ are the parts of $J_\omega$ on which $\alpha_{uu'}$ is smaller,
respectively larger, than at $m$.
\item Let $A$ be an arc of the model arc system of $e=[n,n']$, with edge $[u,u']$ of $\mathcal T$, and let $n_y\in N$
complete $n'-n$ to a basis of the lattice $\{x\in N:\langle u,x\rangle=\langle u',x\rangle=0\}$. Then
$(w_u,w_{u'},n'-n,n_y)$ is a basis of $N$, and the level lines of $y:=\langle\,\cdot\,,n_y\rangle$ in $E_e$ are parallel to
$u'-u$, so $A$ is the graph of a piecewise linear function $\bar\alpha$ over an interval of values of $y$:
$\alpha_{uu'}=\bar\alpha(y)$ on $A$. Put $y_e:=\operatorname{val}_{n'}-\operatorname{val}_n$,
$\alpha_e:=\alpha_{uu'}-\bar\alpha(y)$ and $T_e:=(y_e,\alpha_e,y)$. Near the part of $A$ in the interior of $E_e$, $T_e$
is a piecewise linear homeomorphism onto an open subset of $\R^3$, mapping $\operatorname{int}F^*_n$, $E_e$ and
$\operatorname{int}F^*_{n'}$ into $\{y_e<0\}$, $\{y_e=0\}$ and $\{y_e>0\}$, $A$ onto a segment of the $y$-axis, and the
cells $C^M(u,e)$, $C^M(u',e)$ near $A$ into $\{y_e=0,\ \alpha_e<0\}$ and $\{y_e=0,\ \alpha_e>0\}$.
\end{enumerate}
\end{lemma}

The proof is in Appendix~\ref{app:liftproofs}.

So at every model node and model vertex the side $\alpha_{uu'}>\alpha_{uu'}(m)$ is the side of $u'$, the side of the
tentacles of the divisor $D_{u'}$. The model nodes are the midpoints of pairwise distinct edges $J_P$ of $B_Y$; this
includes the two nodes of a pentagon and the three nodes of an $A_2$ hexagon, which share the edge $[u,u']$ of
$\mathcal T$. We take the $\varrho_q$ of Lemma~\ref{lem:nodemodel}(4) so small that the cubes
$T_q^{-1}([-\varrho_q,\varrho_q]^3)$ are pairwise disjoint and disjoint from fixed neighbourhoods of the other
$\Gamma^M$-vertices.

\begin{lemma}[model cells near the model discriminant]\label{lem:place}
There is $r_S\in(0,\theta_4/4]$, depending only on the configuration, the model positions and $\Upsilon^M$, such that:
\begin{enumerate}
\item at each model node $m_P$, $T_q^{-1}(\{y_1=0,\ \alpha_q\ge0\})\cap T_q^{-1}([-r_S,r_S]^3)$ is the intersection with
that cube of the closure of $C^M(u',ab)\cup C^M(u',P)\cup C^M(u',cd)$, and it is bounded on $\Gamma^M$ by the model arcs of
the legs over $ab$ and $dc$; likewise $\{y_2=0,\ \alpha_q\ge0\}$ with $bc$ and $da$, and $\{\alpha_q\le0\}$ with $u$ in
place of $u'$;
\item near each model vertex $m=m_{\omega,i}$, $\omega=[n_a,n_b,n_c]$, with edge $[u,u']=[u_{i-1},u_i]$ of $\mathcal T$ and
edges $e_1,e_2,e_3$ of $\omega$, the closure of $C^M(u',\omega)\cup\bigcup_kC^M(u',e_k)$ is the union of three closed
sectors, one in each $E_{e_k}$, meeting along the half-edge $C^M(u',\omega)$ of $J_\omega$, the $k$-th bounded by that
half-edge and the model arc $C^M([u,u'],e_k)$; likewise with $u$;
\item near each turning point $C^M(\tau,e)$, $\tau=[u_1,u_2,u_3]$, the closure of $C^M(u_i,e)$ is the sector of $E_e$ between
the model arcs of $[u_i,u_j]$ and $[u_i,u_k]$, of angle less than $\pi$;
\item near each model arc $A$ in the interior of $E_e$, in the chart $T_e$, the closures of $C^M(u',e)$ and $C^M(u,e)$ are
$\{y_e=0,\ \alpha_e\ge0\}$ and $\{y_e=0,\ \alpha_e\le0\}$.
\end{enumerate}
\end{lemma}

The proof is in Appendix~\ref{app:liftproofs}.

\subsubsection*{The transferred cycle}
We apply Definition~\ref{def:transfer} with $\Upsilon=\Upsilon^M$: $\mathfrak h=\Upsilon^M\circ\mathcal L\circ g_C\colon B'_\sigma\to
B_Y$, $\kappa_{\mathfrak h}\colon\Lambda\to\mathfrak h^*(\Lambda^M)^\vee$, the coefficient groups $F^*(s)$ and $\mathcal D^*(E)=\Z d^*_A$,
$d^*_A=\langle\,\cdot\,,n'-n\rangle$ on the arc $A$ of $([u,u'],[n,n'])$. For $z\in Z(\mathfrak K)$ the \emph{transferred
cycle} is
\[
z^M:=\mathfrak h_*z\in Z\bigl(\mathfrak h(\mathfrak K_1);F^*,\mathcal D^*\bigr),
\]
a relative cycle on the triangulation $\mathfrak h(\mathfrak K_1)$ of $B_Y$, in which $\Gamma^M$ is full. By
Proposition~\ref{prop:transfer} its boundary network $b^M:=\partial z^M$ has the integers of $\partial z$, and
$\rho^*(z^M)=\rho(z)$; we write $\rho$ for $\rho^*$ and $\mathrm{Bal}(\Gamma^M)$ for the group of balanced networks on
$\Gamma^M$, the $1$-cycles of $C_*(\Gamma^M;\mathcal D^*)$. By Proposition~\ref{prop:dict}(2),(3) the stalks
$\operatorname{Inv}^*$ are: on an arc, $\{v:v(u'-u)=0\}$, of rank two, with basis $d^*_A$ and $g^*_A:=\langle\,\cdot\,,n_y
\rangle$ (Lemma~\ref{lem:crossing}); at a model node, the same group, with basis $\langle\,\cdot\,,b-a\rangle$,
$\langle\,\cdot\,,c-b\rangle$; at a model vertex, of positive type, a group of rank two containing the covectors $d^*_A$ of
its three arcs; at a turning point $C^M(\tau,e)$, of negative type, $\Z d^*_e$, $d^*_e:=\langle\,\cdot\,,n'-n\rangle$ for
$e=[n,n']$, the common covector of its three arcs up to sign.

Lemmas~\ref{lem:nbhd} and~\ref{lem:normalform} hold on $B_Y$, for $\Gamma^M$ and the coefficients $F^*$, $\mathcal D^*$,
on every triangulation in which $\Gamma^M$ is full: their proofs use only a triangulated three-manifold, a full graph and
the flat sections of a local system (Remark~\ref{rem:nbhdscope}). Equivalently, for a triangulation $\mathfrak K^*$ of $B_Y$
refining $\mathfrak h(\mathfrak K_1)$ the map $\mathfrak h^{-1}$ is linear on the simplices of $\mathfrak K^*$, and Lemma~\ref{lem:transport}
applied to $(\mathfrak h^{-1},\kappa_{\mathfrak h}^{-1})$ identifies the chain complexes, coefficient systems and stalks on $\mathfrak K^*$ with those
on the triangulation $\mathfrak h^{-1}(\mathfrak K^*)$ of $B'_\sigma$, in which $\Gamma$ is full.

\begin{lemma}[balanced leg data]\label{lem:unitgerms}
At a $\Gamma^M$-vertex $p$ let $\mathrm{Bal}_p$ be the lattice of \emph{balanced leg data}: the vectors $(c_A)$ of coefficients on
the arcs $A$ at $p$, oriented away from $p$, with $\sum_Ac_Ad^*_A=0$ in $\operatorname{Inv}^*(p)$. Then: at a model vertex,
with the $d^*_A$ oriented so that $d^*_1+d^*_2+d^*_3=0$, $\mathrm{Bal}_p=\Z(1,1,1)$; at a turning point, with the $d^*_A$ oriented
equal, $\mathrm{Bal}_p=\{c_{12}+c_{13}+c_{23}=0\}$; at a model node, with $d^*_{ab}=d^*_{dc}=\langle\,\cdot\,,b-a\rangle$ and
$d^*_{bc}=d^*_{da}=\langle\,\cdot\,,c-b\rangle$, $\mathrm{Bal}_p=\{c_{ab}=-c_{dc},\ c_{bc}=-c_{da}\}$. A balanced network
restricts at $p$ to an element of $\mathrm{Bal}_p$.
\end{lemma}

\begin{proof}
At a model vertex the three $d^*_A$ are $\pm\langle\,\cdot\,,e_k\rangle$ for the three edges $e_k$ of a unimodular triangle,
which span a lattice of rank two and satisfy one relation; at a turning point they are equal up to sign; at a node
$\langle\,\cdot\,,b-a\rangle$ and $\langle\,\cdot\,,c-b\rangle$ are independent. The last statement is the balance of a balanced network
at $p$, a relation in $F^*(p)=\operatorname{Inv}^*(p)$. This is Lemma~\ref{lem:legdata}, transported by $\mathfrak h$ and
$\kappa_{\mathfrak h}$, with the
vertex types exchanged.
\end{proof}

\subsubsection*{Model germs}
All germs below are fixed once, after the configuration, the model positions and $\Upsilon^M$, and before $z$ and $t$.
They lie in $N_0$, the $r_0$-neighbourhood of $\Gamma^M$, where $r_0\le r_S/4$ is so small that $N_0$ lies in the union of
the cubes $T_q^{-1}([-r_S,r_S]^3)$, the neighbourhoods of the model vertices and turning points of
Lemma~\ref{lem:place}(2),(3), and the domains of the charts $T_e$. Each germ is a piecewise linear surface with boundary
on $\Gamma^M\cup\partial N_0$, carrying a coefficient that is a flat section of $(\Lambda^M)^\vee$ over it. The
\emph{away-oriented leg data} of a germ at a $\Gamma^M$-vertex $p$ are the coefficients $c_A$ of its boundary
$\sum_Ac_Ad^*_A[A]$ on the arcs $A$ at $p$, each oriented away from $p$.

\begin{enumerate}
\item[(M1)] \emph{Node halves.} At a model node $m_P$ the four \emph{standard halves} are
$T_q^{-1}(\{y_1=0,\ \pm\alpha_q\ge0\})$ (the \emph{$b-a$ passage}, along the legs over $ab$ and $dc$) and
$T_q^{-1}(\{y_2=0,\ \pm\alpha_q\ge0\})$ (the \emph{$c-b$ passage}, along $bc$ and $da$). We use the side $\alpha_q\ge0$.
The \emph{unit germ} $g_{ab}$ is the $b-a$ half with the orientation of the frame $(y_2,\alpha_q)$ and the covector
$-\langle\,\cdot\,,b-a\rangle$; the unit germ $g_{bc}$ is the $c-b$ half with the orientation of the frame
$(y_1,\alpha_q)$ and the covector $+\langle\,\cdot\,,c-b\rangle$. The boundary of $g_{ab}$ runs from the leg over $ab$ to the leg over $dc$, and that of
$g_{bc}$ from the leg over $da$ to the leg over $bc$; so their away-oriented leg data are $(c_{ab},c_{bc})=(1,0)$ and
$(0,1)$, and $\rho_q(g_{ab})=1$, $\rho_q(g_{bc})=-1$. Both covectors lie in $\operatorname{Inv}^*(m_P)$.
\item[(M2)] \emph{Tripods.} At a model vertex $m=m_{\omega,i}$ the unit germ is the union of the three sectors of
Lemma~\ref{lem:place}(2) on the side of $u'$. Orient the edges $e_k$ of $\omega$ cyclically, so that $e_1+e_2+e_3=0$, orient
the sector in $E_{e_k}$ so that its boundary leaves $m$ along its model arc (and hence arrives at $m$ along the half-edge
$C^M(u',\omega)$), and give it the covector $\langle\,\cdot\,,e_k\rangle$, which is $d^*_A$ for its arc. The three sectors
induce the same orientation on the half-edge, their covectors sum to zero there, and the away-oriented leg data are
$(1,1,1)$.
\item[(M3)] \emph{Turning germs.} At a turning point $C^M(\tau,e)$, $\tau=[u_1,u_2,u_3]$, $e=[n,n']$, let $A_{ij}$ be the arc
of $[u_i,u_j]$ and $C(u_i)$ the sector of Lemma~\ref{lem:place}(3) between $A_{ij}$ and $A_{ik}$. Fix an orientation $o_e$
of $E_e$. With every sector oriented by $o_e$ and carrying $d^*_e$, the leg data $(c_{12},c_{13},c_{23})$ of the three
sectors sum to zero, and each sector has data $\pm1$ on its two arcs and $0$ on the third. We label $\tau$ so that
$C(u_1)=(1,-1,0)$, $C(u_2)=(-1,0,1)$, $C(u_3)=(0,1,-1)$. The \emph{turning germs} are $\mathrm{Turn}(u_1):=C(u_1)$ and
$\mathrm{Turn}(u_3):=C(u_3)$ with orientation $o_e$ and $\mathrm{Turn}(u_2):=C(u_2)$ with orientation $-o_e$, all with
covector $d^*_e$; their leg data are $(1,-1,0)$, $(1,0,-1)$, $(0,1,-1)$. The chain $\mathrm{Turn}(u_1)-\mathrm{Turn}(u_2)+
\mathrm{Turn}(u_3)$, the full disc of $E_e$ about the vertex with orientation $o_e$, has zero leg data.
\item[(M4)] \emph{Exit sheets and reference sheets.} Orient each arc $A$ from one end $p_0$ to the other end $p_1$ and fix,
in the part of $A$ in $\Omega_e$, five points in this order:
$y^*_{A,0},\ y^0_A,\ \mathfrak x_A,\ y^1_A,\ y^*_{A,1}$, and take $r_0$ at most a quarter of their mutual distances, of
their distances from $A\setminus\Omega_e$ and of their distances from the ends $p_0,p_1$ of $A$, so that the stretches of
length at most $r_0$ about $y^*_{A,p}$ below lie in $\Omega_e$ and away from the balls $B^0_{p_0}$, $B^0_{p_1}$ of
Section~\ref{ssec:lift-collar}, whose radius is at most $r_0$. A unit germ at an end $p$ of $A$ that uses $A$ continues along $A$
as its \emph{exit sheet}: in the chart $T_e$, the half-plane $\{y_e=0,\ \alpha_e\ge0\}$ (the closure of $C^M(u',e)$) or
$\{y_e=0,\ \alpha_e\le0\}$ (that of $C^M(u,e)$), according to its side at $p$, from $p$ to $y^0_A$ if $p=p_0$ and to
$y^1_A$ if $p=p_1$. The \emph{reference side} of every arc is $\alpha_e\ge0$. Fix a convex polygon $\mathsf P\subset\R^2$,
symmetric about $0$, with gauge $\varrho_{\mathsf P}$. An exit sheet on the side $\alpha_e\le0$ is \emph{rotated} to the
reference side at $y^*_{A,p}$: in the chart it is the half-plane $\alpha_e\le0$ up to $y^*_{A,p}$, the normal half-disc
$\{y=y^*_{A,p},\ y_e\ge0,\ \varrho_{\mathsf P}(y_e,\alpha_e)\le r_0\}$ there, and the half-plane $\alpha_e\ge0$ after it.
All half-planes are truncated at $\varrho_{\mathsf P}\le r_0$. The \emph{reference sheet} of $A$ with coefficient
$c_Ad^*_A$ is the half-plane $\alpha_e\ge0$ over $[y^0_A,y^1_A]$, oriented so that its boundary on $A$ runs in the direction
of $A$.
\item[(M5)] \emph{Transverse discs.} At $\mathfrak x_A$, the normal disc $\mathfrak N_A:=\{y=y(\mathfrak x_A),\
y_e^2+\alpha_e^2\le r_0^2\}$ in the chart $T_e$, cooriented by $A$: (orientation of $\mathfrak N_A$, direction of $A$) is
the orientation of $B_Y$.
\end{enumerate}

At an end $p$ of $A$ the exit sheets on a given side come from one germ type: at a turning point the side $u_i$ of the
arc of $[u_i,u_j]$ belongs to $\mathrm{Turn}(u_i)$, and at a model node or vertex each leg is used by one germ type, on
the side $u'$. An arc joining two model nodes receives one exit sheet from each end, cut at $y^0_A$ and $y^1_A$. The rotation through the half-plane $y_e>0$ rather than through $y_e<0$ is a convention; the other choice changes
the chain by a transverse disc with coefficient in $\Z d^*_A$.

By Lemma~\ref{lem:unitgerms} the unit germs at $p$ form a $\Z$-basis of $\mathrm{Bal}_p$: $g_{ab},g_{bc}$ at a model node, the
tripod at a model vertex, $\mathrm{Turn}(u_1),\mathrm{Turn}(u_2)$ at a turning point. For $b\in\mathrm{Bal}(\Gamma^M)$ write
$b_p=\sum_gm_g(b)\,\partial g$ at every $\Gamma^M$-vertex $p$, with integers $m_g(b)$ ($g$ running over the
unit germs of the basis at $p$).

\begin{definition}[the standard chain]\label{def:std}
For $b\in\mathrm{Bal}(\Gamma^M)$ the \emph{standard chain} $\mathrm{Std}^M(b)$ is the sum over all $\Gamma^M$-vertices $p$ of
$\sum_gm_g(b)\,g$, each unit germ with its exit sheets, cut transversally at $y^0_A$ or $y^1_A$, plus, on every arc
$A$, the reference sheet with covector $c_Ad^*_A$ between $y^0_A$ and $y^1_A$, where $c_Ad^*_A[A]$ is the part of $b$
on $A$.
\end{definition}

Along $A$ the exit sheets arriving from each end carry, on the reference side, covectors summing to $c_Ad^*_A$ with the
orientation of the reference sheet, by the balance of $b$ along the arc; so the transverse cut segments cancel, and
$\mathrm{Std}^M(b)$ is a relative cycle of $(N_0,\partial N_0\cup\Gamma^M)$ with boundary network $b$. It is linear in
$b$.

Where sheets of $\mathrm{Std}^M(b)$ coincide as sets, namely sheets on the reference side of one arc, their
coefficients add.

\subsubsection*{The normal form on \texorpdfstring{$B_Y$}{BY}}
Let $z^M=\mathfrak h_*z$ and $b^M=\partial z^M$. Choose a triangulation $\mathfrak K_2^0$ of $B_Y$ refining $\mathfrak h(\mathfrak K_1)$, in which
$\Gamma^M$, the germs (M1)--(M5) and their junction segments are subcomplexes, and which is so fine near $\Gamma^M$ that
the simplicial neighbourhoods below lie in the interior of $N_0$. Let $\mathfrak K_2$ be a derived subdivision of
$\mathfrak K_2^0$, in which $\Gamma^M$ is full and the germs are subcomplexes; choose its derived points off the points
$\mathfrak x_A$, so that $\mathfrak x_A$ lies in the interior of an edge $E_A$ of $\mathfrak K_2$ in $A$; and let
$\mathfrak K_2'$ be a derived subdivision of $\mathfrak K_2$ whose new vertex on $E_A$ is $\mathfrak x_A$. Put
$U:=N(\Gamma^M,\mathfrak K'_2)$, the union of the closed simplices meeting $\Gamma^M$. Then $\operatorname{sd}z^M$ lies on
$\mathfrak K_2'$, and $\mathrm{Std}^M(b^M)\cap U$ is a relative cycle of $(U,\partial U\cup\Gamma^M)$ on $\mathfrak K'_2$ with
boundary network $b^M$: its coefficients lie in the stalks, since near a $\Gamma^M$-vertex only the germs of that vertex
occur and along an arc only covectors in $\operatorname{Inv}^*(A)$; and an edge of $U$ not in $\partial U$ meets
$\Gamma^M$, so that all $2$-simplices at it lie in $U$ and intersecting with $U$ creates boundary only on $\partial U$. By
Lemma~\ref{lem:normalform}, with the edge $E_A$ chosen in each arc $A$,
\begin{equation}\label{eq:nf1}
\operatorname{sd}z^M-\partial w=\mathrm{Std}^M(b^M)\cap U+\sum_Av_A\,\bar D(E_A)+y,\qquad v_A\in\operatorname{Inv}^*(A),
\end{equation}
with $w$ a $3$-chain in $U$, $v_A\bar D(E_A)$ the transverse disc at $E_A$ in $\mathfrak K'_2$ (Definition~\ref{def:transverse}), centred at $\mathfrak x_A$, and $y$
supported in $B_Y\setminus\operatorname{int}U$. Of Lemma~\ref{lem:nbhd} only the surjectivity in part~(1) enters, applied to the chain
$\operatorname{sd}z^M|_U-\mathrm{Std}^M(b^M)\cap U$, which has zero boundary on $\Gamma^M$.

We then make two local replacements. A \emph{local replacement} replaces a chain $C$ by a chain $C'$ with the same
boundary on $\Gamma^M$ and adds a chain $R$ supported off $\Gamma^M$ with $\partial R=\partial C-\partial C'$ off $\Gamma^M$;
it preserves the boundary network.
\begin{itemize}
\item \emph{Straightening the transverse discs.} Let $N_{\varepsilon_A}(\mathfrak x_A)$ be a piecewise linear regular
neighbourhood of $\mathfrak x_A$ of diameter at most $\varepsilon_A$, for instance a closed star in a derived subdivision in
which $\bar D(E_A)$ and $\Gamma^M$ are subcomplexes. By regular-neighbourhood theory \cite[Ch.~3]{RS} the triple
$(N_{\varepsilon_A},N_{\varepsilon_A}\cap\bar D(E_A),N_{\varepsilon_A}\cap\Gamma^M)$ is unknotted (a ball, a proper disc and a
proper arc meeting it once transversally), so $\partial N_{\varepsilon_A}$ minus the two points of $\Gamma^M$ is an open
annulus in which $\partial N_{\varepsilon_A}\cap\bar D(E_A)$ is essential. Let $\mathfrak N_A^{\varepsilon'_A}$ be the normal
disc of (M5) of radius $\varepsilon'_A\le r_0$, so small that its solid cylinder $T_e^{-1}(\{y_e^2+\alpha_e^2\le
\varepsilon'^2_A,\ |y-y(\mathfrak x_A)|\le\varepsilon'_A\})$ lies in $\operatorname{int}N_{\varepsilon_A}$. Replace
$v_A(\bar D(E_A)\cap N_{\varepsilon_A})$ by $v_A\mathfrak N^{\varepsilon'_A}_A$ plus $v_A$ times an embedded annulus in
$N_{\varepsilon_A}\setminus(\Gamma^M\cup\text{cylinder})$ between the two essential circles; $v_A$ extends over it because it
lies in a punctured neighbourhood of an arc point.
\item \emph{Smoothing the rotations.} On a stretch $\mathfrak a_{A,p}\ni y^*_{A,p}$ of length at most $r_0$ replace every
rotated exit sheet by the union of the rays $\{(\varrho\cos\Theta(y),\varrho\sin\Theta(y),y):\varrho\ge0\}$ in the
coordinates $(y_e,\alpha_e,y)$, truncated at $\varrho_{\mathsf P}\le\varrho_0$, with $\Theta$ smooth and monotone,
$\Theta=-\pi/2$ near the start and $\Theta=\pi/2$ near the end of the stretch, passing through $\Theta=0$; add the strip on
the polygonal cylinder $\{\varrho_{\mathsf P}=\varrho_0\}$ over the stretch between the two outer curves, closed up at its
ends.
\end{itemize}
The sizes $\varepsilon_A$ and $\varrho_0$ are chosen so that the neighbourhoods $N_{\varepsilon_A}(\mathfrak x_A)$ and the
solid cylinders $\{\varrho_{\mathsf P}\le\varrho_0\}$ over the stretches lie in $\operatorname{int}U$, and each meets no
summand other than the sheets along its own arc.

\begin{lemma}[normal form of the transferred cycle]\label{lem:nfY}
The result is a relative cycle
\begin{equation}\label{eq:nf2}
z^\natural=\mathrm{Std}^M(b^M)+\sum_Av_A\,\mathfrak N_A+y^\natural,\qquad v_A\in\operatorname{Inv}^*(A),
\end{equation}
with boundary network $b^M$ and $\rho(z^\natural)=\rho(z)$, where $\mathrm{Std}^M(b^M)$ now has smooth rotations and is
truncated at $\partial U$, $\mathfrak N_A$ is $\mathfrak N^{\varepsilon'_A}_A$ continued by the annulus and by $\bar D(E_A)$
outside $N_{\varepsilon_A}(\mathfrak x_A)$, and $y^\natural$ is $y$ plus the strips of the smoothing. Let $r_\natural>0$ be
the distance, in a fixed metric on $B_Y$, from $\Gamma^M$ to the union of $\operatorname{supp}y^\natural$ and of the parts
of the transverse discs outside the normal discs $\mathfrak N^{\varepsilon'_A}_A$. Within distance $r_\natural$ of
$\Gamma^M$, $z^\natural$ is exactly $\mathrm{Std}^M(b^M)+\sum_Av_A\mathfrak N^{\varepsilon'_A}_A$: germs in the positions of
Lemma~\ref{lem:place} and normal discs at the points $\mathfrak x_A$.
\end{lemma}

\begin{proof}
The local replacements preserve the boundary network. The strips of the smoothing lie at distance at least
$\varrho_0/C_{\mathsf P}$ from $\Gamma^M$, $C_{\mathsf P}$ the equivalence constant of $\varrho_{\mathsf P}$ and the
Euclidean norm; the annuli lie in compact subsets of $N_{\varepsilon_A}\setminus\Gamma^M$; $y$ lies off
$\operatorname{int}U$; and the transverse discs meet $\Gamma^M$ only at their centres. So $r_\natural>0$, with
$r_\natural\le\min(\varepsilon'_A,\varrho_0/C_{\mathsf P})$.
\end{proof}

The homology class of $z$ plays no role: what follows needs only a relative cycle with the boundary network of $z$, and
Lemma~\ref{lem:nfY} provides one of the required shape.

\emph{Summands.} Near $\Gamma^M$ the normal form~\eqref{eq:nf2} is a finite sum of surfaces with integer
coefficients: for each $\Gamma^M$-vertex $p$ and each unit germ $g$ at $p$, the surface of $g$ (its component at $p$ in
$N_0$ together with its exit sheets, rotated and smoothed where required, up to $y^0_A$ or $y^1_A$), with coefficient
$m_g(b^M)$ times the covector of $g$; the reference sheets, with coefficients $c_Ad^*_A$; and the transverse discs
$v_A\mathfrak N_A$. Where surfaces of this list coincide as sets, which happens only for sheets on the reference side of
one arc, we add their coefficients. A \emph{summand} is a maximal connected part of one of these surfaces, or of a
union of coincident sheets, on which the total coefficient is constant, except that the three sectors of a tripod, with
their exit strips, form one summand: its half-edge $C^M(u',\omega)$ is an interior branching curve, along which the
three coefficients $m\langle\,\cdot\,,e_k\rangle$ sum to zero with the orientations of (M2). In particular the part of a unit germ at its
vertex and its exit strips, including the start of a rotation, belong to one summand up to the first place where
another sheet joins them: the transverse segments where a germ continues as its own exit strip carry the same
coefficient on both sides and are interior to one summand. Two summands are \emph{adjacent} along a
\emph{junction segment}, a curve other than the half-edge of a tripod across which the total coefficient on one side of
the arc changes. Junction segments occur only where a rotated exit strip reaches a reference side that already carries a
strip (inside the polygonal cylinder at the far end of the stretch; outside it, where the piecewise linear normal
half-disc of the rotation at $y^*_{A,p}$ meets the reference-side strip). Where a rotated strip reaches an empty
reference side, it simply continues there, and where the exit strips on the reference side meet the reference sheet, at
$y^0_A$ and $y^1_A$, they already total $c_Ad^*_A$ with the orientation of the reference sheet (by the balance of $b^M$
along $A$); in both cases the total coefficient does not change and the curve is interior to one summand. Along a
junction segment every
adjacent coefficient is a multiple of $d^*_A$, and the multiples sum to zero with the orientation signs, because
$z^\natural$ is a relative cycle with boundary only on $\Gamma^M$. Inside the stretch, after $y^*_{A,p}$, the reference-side
sheet carries one coefficient inside the cylinder $\{\varrho_{\mathsf P}<\varrho_0\}$, where the smoothed strip is still
rotating, and another outside it, where the piecewise linear rotated sheet has already arrived; they meet the correction
strip of the smoothing along the seam $\{\varrho_{\mathsf P}=\varrho_0\}$, at distance at least $\varrho_0/C_{\mathsf P}\ge
r_\natural$ from $\Gamma^M$. Summands may also meet transversally (a transverse disc and a sheet; the two node halves
along $C^M(u',P)$; exit sheets of different germs meeting only on $\Gamma^M$); such meetings are not junctions.

\subsection{Pieces along legs, at vertices and at turning points}\label{ssec:lift-legs}

We lift the boundary germs in the binomial and pair-of-pants models of the localized hypersurface.
The positive real section fixes compatible base points for the torus orbits, and Lemma~\ref{lem:windows} places these orbits in regions where the equations are exact.

\subsubsection*{The positive section}
\begin{lemma}[the positive locus]\label{lem:possec}
For $t$ small and every real $c$ with the sign pattern:
\begin{enumerate}
\item the positive real locus $\Sigma^+:=Y_t\cap(\R_{>0})^4$ is a smooth strictly convex three-sphere in the torus
(compare \cite[\S3.1]{AGIS20});
\item for $s\in[0,1]$ the positive locus $\Sigma^+_s$ of $W_s$ is a smooth closed three-manifold, the $\Sigma^+_s$ form a
smooth family, and $\operatorname{Log}(\Sigma^+_0)/L$ is a star-shaped sphere within $O(1/L)$ of $B_Y$. Radial projection
therefore defines a section $\mathfrak s^\vee\colon B_Y\to\Sigma^+_W:=\Sigma^+_0\subset W^*$.
\end{enumerate}
\end{lemma}

\begin{proof}
(1) On the positive part of the torus, $F_t/(\gamma_0x^0)=1-\sum_{n\ne0}|\gamma_n/\gamma_0|e^{\langle w,n\rangle}$ with
$w=\operatorname{Log}x$. The sum is strictly convex in $w$, because the $n\ne0$ span $N_\R$, and proper, because $0$ is
interior to their convex hull; its value at $w=0$ is $O(t^{k_0+\min h_C})<1$. So $\{\sum=1\}$ is a smooth strictly convex
sphere.

(2) Complex conjugation preserves $F^{(s)}$, since the coefficients are real and the weights depend only on moduli; so
at a real point $dF^{(s)}$ maps real tangent vectors to $\R$ and imaginary ones to $i\R$, and real rank two
(Lemma~\ref{lem:W}(1)) forces $dF^{(s)}|_{\R^4}\ne0$. Hence $\Sigma^+_s$ is a regular level set, and the family is smooth.
It does not meet the toric boundary: on the nonnegative part of a divisor the origin monomial vanishes and the others are
nonpositive, the maximal one having weight one. Along a ray $w=\lambda w_0$, at a zero of $F^{(0)}$ the origin monomial
lies between $\mathfrak m$ and $|\mathcal A|\mathfrak m$, $\mathfrak m$ the largest monomial; so the deficit of the origin is
$O(1/L)$, $w/L$ lies within $O(1/L)$ of $B_Y$, and $|w|=O(L)$. A monomial $n\ne0$ of positive weight has
$d_n\le\tilde d_n<\theta_1/2\le1/2$, so $|\gamma_nx^n|\ge t^{1/2}\mathfrak m$, that is $\langle w,n\rangle\ge
L(H(n)-\tfrac12)-O(1)\ge\tfrac L2-O(1)$, since $H(n)\ge1$ by~\eqref{eq:lambdagap}. The function
$\mathfrak S(\lambda):=\sum_{n\ne0}\psi_n|\gamma_n/\gamma_0|e^{\langle w,n\rangle}$ equals $1$ at a crossing; with the weights
held constant its derivative with respect to $\log\lambda$ is at least $L/2-O(1)$ there, and the derivatives of the weights
are $O(1)$ (Lemma~\ref{lem:refined}(2)) and multiply monomials in transition, of total size $O(t^{\theta_1/4})$ relative to
the origin. So each ray crosses $\Sigma^+_0$ exactly once, transversally.
\end{proof}

For a primitive covector $v$ in the stalk of $(\Lambda^M)^\vee$ at a point of $B_Y\setminus\Gamma^M$, the quotient of the
fibre torus $\Lambda^M\otimes U(1)$ by $\mathbb T(v):=\ker v\otimes U(1)$ is the circle of the character $\chi^{\tilde v}$
for any lift $\tilde v\in N$ of $v$, and its phase is $0$ on $\Sigma^+_W$. So the union of the $\mathbb T(v)$-orbits
through the points of $\mathfrak s^\vee$ over a surface with coefficient $v$ is contained in $\{\arg\chi^{\tilde v}=0\}$.

\subsubsection*{Where the localized mirror is given by exact models}
We use Lemma~\ref{lem:windows}: on the sets listed there the weights of the localized mirror are exactly $0$ or $1$. For
node parallelograms, item (1$'$) covers the stretches of the model arcs from a model node to $\Omega_e$, along the leaf
paths, and for triangles those from a model vertex.

\subsubsection*{Normalised coordinates along the model arcs}
Let $A$ be an arc of the model arc system of $e=[n,n']$, with edge $[u,u']$ of $\mathcal T$ and $n_y$ as in
Lemma~\ref{lem:tropcharts}(2). On the chart $U_{u,u'}$ we use $Z=\chi^{w_u}$, $Z'=\chi^{w_{u'}}$, $\chi_1=\chi^{n'-n}$ and
$\chi_2=\chi^{n_y}$. The \emph{normalised coordinates} along $A$ are $\hat Z=Z/\nu_Z$ and $\hat Z'=Z'/\nu_{Z'}$, where
$\nu_Z,\nu_{Z'}>0$ are smooth functions of $\log|\chi_1|$ and $\log|\chi_2|$ such that $|\gamma_0x^0|=|\gamma_nx^n|\,|\hat
Z\hat Z'|$, and $\nu_Z/\nu_{Z'}$ is a function of $\log|\chi_2|$ alone with $\frac1L\log(\nu_Z/\nu_{Z'})=\bar\alpha(\log
|\chi_2|/L)+O(1/L)$; where the torus coordinates are within $2\theta_4$ of those of a model node they are the constants of
Definition~\ref{def:morsecoords}, with the reference monomial $a$ in place of $n$, and near a model vertex they are
constants with $|\hat Z|=|\hat Z'|$ at the vertex; on these sets $\bar\alpha$ is constant, so the two requirements agree.
Put $\hat\mu:=|\hat Z|^2-|\hat Z'|^2$ and
\[
\mathcal S_G:=\sum_m\psi_m\,|\gamma_mx^m|/|\gamma_nx^n|,
\]
the sum over the lattice points $m$ of the face $[u,u']^\vee$ of $\Delta$, the monomials free of $Z$ and $Z'$. Changing
the reference monomial $n$ to $n''$, with the same ratio $\nu_Z/\nu_{Z'}$, multiplies $|\hat Z|^2$, $|\hat Z'|^2$ and
$\mathcal S_G$ by the same function of $(\chi_1,\chi_2)$. Where only $0,n,n'$ have positive weight, $W^*$ is the \emph{leg
model}
\begin{equation}\label{eq:legmodel}
1+\kappa_e\chi_1+\mathfrak b\,ZZ'=0,\qquad\kappa_e=\gamma_{n'}/\gamma_n>0,\qquad|\mathfrak b|\nu_Z\nu_{Z'}=1,
\end{equation}
with $\mathfrak b$ a monomial in the torus coordinates, and $\mathcal S_G=1+|\kappa_e\chi_1|$. Its critical locus is
$\{Z=Z'=0,\ \chi_1=-1/\kappa_e\}$, and the rank-two torus $\mathbb T_d:=(\{n,n'\}^\perp\cap M)\otimes U(1)$, the
\emph{exact torus} of the leg, acts on it; it contains $S^1_\alpha$, and $\{n,n'\}^\perp\cap M=\ker d^*_A$ is the sublattice
of $\Lambda^M$ fixed by the monodromy around the arc.

The \emph{leg model base} is the $(x_1,x_2)$-plane, $x_1=\log|\kappa_e\chi_1|$ and $x_2=\hat\mu$, times the leg
coordinate $y$. As at a node, $n$ is the larger of $n,n'$ where $x_1<0$ and $n'$ where $x_1>0$, and $x_2<0$ (resp.\
$x_2>0$) is the side of $D_u$ (resp.\ $D_{u'}$). At a point of the torus, $T_e(\operatorname{Log}x/L)=(x_1,\
\log|\hat Z/\hat Z'|,\ \log|\chi_2|)/L+O(1/L)$. The \emph{model map}
\[
E_t(y_e,\alpha_e):=\bigl(Ly_e,\ (1+e^{Ly_e})\sinh(L\alpha_e)\bigr)
\]
is a smooth homeomorphism of $\R^2$ with $E_t(y_e,\alpha_e)=L(y_e,2\alpha_e)+O(L^2(y_e^2+\alpha_e^2))$ near $0$; it
takes the tropical coordinates $(y_e,\alpha_e)$ to the coordinates $(x_1,x_2)$ of the leg model base, in the sense of
Lemma~\ref{lem:interface}.

\subsubsection*{The pieces}
\begin{lemma}[pieces along the arcs, at vertices and at turning points]\label{lem:pieces}
Let $0<\varsigma\le\varsigma_{\max}$ for the constant $\varsigma_{\max}\le\theta_2/4$ of the proof, which depends only on the
configuration and the model arc systems, and $t$ small.
\begin{enumerate}
\item \emph{Arcs.} Let $\mathfrak a$ be a compact stretch of a model arc $A$ of $e=[n,n']$, contained in $\Omega_e$, and
$\Theta\colon\mathfrak a\to\R$ smooth. For the sheet $\{(\varrho\cos\Theta(y),\varrho\sin\Theta(y),y):0\le\varrho\le\varrho_0,
\ y\in\mathfrak a\}$ in the chart $T_e$, $\varrho_0\ge\varsigma$, put
\[
\Pi_{\mathfrak a,\varsigma}:=\mathbb T_d\cdot\bigl\{\mathbf x\bigl(E_t(\varrho\cos\Theta(y),\varrho\sin\Theta(y)),\,y\bigr):
0\le\varrho\le\varsigma,\ y\in\mathfrak a\bigr\},
\]
where $\mathbf x(x_1,x_2,y)$ is the point of $W^*$ with $\log|\chi_2|=Ly$, $\arg\chi_2=0$, $\chi_1=-e^{x_1}/\kappa_e$,
$\hat\mu=x_2$ and $\hat Z\ge0$. Then $\Pi_{\mathfrak a,\varsigma}$ is a smooth four-manifold fibred over the sheet, with
fibres $\mathbb T_d$ collapsing to circles over the arc and no boundary over the arc, and on it the phase of
$\chi^{n'-n}$ is $\pi$. Where $\Theta=\pi/2$ it is the divisor piece $\{Z'=0,\ \chi_1=-1/\kappa_e,\ |\hat Z|^2\le
2\sinh(L\varsigma)\}$, and where $\Theta=-\pi/2$ the corresponding part of $W^*\cap D_u$.

Along a part of a model arc on which the sheet is the closure of $C^M(u',e)$ (resp.\ $C^M(u,e)$), the piece is the
\emph{divisor piece of depth $\varsigma$}: the set of points $x$ of $W^*\cap D_{u'}$ over the stretch and on the branch of
the arc with $|\hat Z(x)|^2\le\mathcal S_G(x)\sinh(L\varsigma)$ (resp.\ the same with $D_u$ and $\hat Z'$). Here the
branch is the curve $\{g_e=0\}$, $g_e$ the weighted sum of the monomials of $[u,u']^\vee$ divided by $\gamma_nx^n$, on
which the two monomials of $e$ are the largest; on the node neighbourhood it is $\mathscr C_P$, and where the torus
coordinates are within $\theta_4$ of those of the node point the divisor piece is the node piece $\Pi^\pm_q$ over
$\mathscr C_P$ with $\bar\mu=\mathcal S_G\sinh(L\varsigma)$. Divisor pieces defined near different model points agree
where they overlap, and they agree with $\Pi_{\mathfrak a,\varsigma}$ where $\Theta=\pm\pi/2$. A divisor piece is a smooth
four-manifold, fibred over the sheet with fibre $S^1_\alpha$ times the circle of the branch, collapsing to that circle
over the arc; on it the phase of $\chi^{n'-n}$ is within $O(\eta_0^{1/2})$ of $\pi$. It is exact: by
Lemma~\ref{lem:windows}(2) on the node neighbourhood, by (1$'$) along the leaf paths, by (1) along the arc in $\Omega_e$
and at the turning points, and by (3) near a model vertex.
\item \emph{Model vertices.} Over $m=m_{\omega,i}$, $\omega=[n_a,n_b,n_c]$, with edge $[u,u']$ of $\mathcal T$,
$W^*=\{\mathfrak bZZ'+\hat f(\chi_1,\chi_2)=0\}$ with $\hat f$ the weighted trinomial of $\omega$. The \emph{vertex piece}
$\Pi_v$ is the set of points of $W^*\cap D_{u'}=\{Z'=0,\ \hat f=0\}$ whose torus coordinates are within $\theta_4$ of those
of $m$ and with $|\hat Z|^2\le\mathcal S_G\sinh(L\varsigma)$: a smooth pair of pants times a disc, as in
\cite[\S7.5]{CBM}. Over the three sectors its fibres are $S^1_\alpha$ times the circle of the tentacle; over the half-edge
$C^M(u',\omega)$ its fibre is $S^1_\alpha$ times the truncated pair of pants, which is the filling $\mathcal Q=\ell\times
\mathcal Q'$ of \cite[\S7.7, Steps~2--3]{CBM} with $\ell=S^1_\alpha$. Along the tentacles it agrees with the divisor pieces
of the three arcs, and its boundary consists of the outer part $\{|\hat Z|^2=\mathcal S_G\sinh(L\varsigma)\}$ and the
ends of the tentacles, where it continues as those divisor pieces.
\item \emph{Turning points.} Over $C^M(\tau,e)$, $\tau=[u_1,u_2,u_3]$, $W^*=\{\mathfrak b\,z_1z_2z_3+\gamma_n(1+\kappa_e
\chi_1)=0\}$ in the Cox coordinates $z_i$ of the $u_i$. The \emph{turning piece} $\Pi_{\rm t}(\mathrm{Turn}(u_1))$ is the
union, near the vertex, of the divisor pieces of depth $\varsigma$ of the arcs of $[u_1,u_2]$ and $[u_1,u_3]$ in
$W^*\cap D_{u_1}=\{z_1=0,\ \chi_1=-1/\kappa_e\}\cong\C^2_{z_2,z_3}$; it is swept by the exact torus rotating $z_2,z_3$, and
near each of the two arcs it is the divisor piece of the common ray $u_1$. It is the piece of \cite[Lemma~7.6]{CBM}. The
pieces of $\mathrm{Turn}(u_2)$, $\mathrm{Turn}(u_3)$ lie in $D_{u_2}$, $D_{u_3}$.
\end{enumerate}
\end{lemma}

The proof is in Appendix~\ref{app:liftproofs}.

\subsection{Node pieces with multiplicity}\label{ssec:lift-node}

The two elementary passages through a node lift to chains whose boundary is the oriented vanishing sphere.
Lemma~\ref{lem:SI} determines the sign; arbitrary integral multiplicities are then handled by addition of chains.

Let $P=abcd$ be a node parallelogram and $q$ its node. We use the node neighbourhood $\mathcal N_P$, the circle
$S^1_\alpha$, the node model and the node pieces $\Pi^\pm_q$ of Section~\ref{ssec:nodepiece}, with the truncation parameter
$\theta'=\theta_4$ of Definition~\ref{def:nodepiece} and, for $0<\varsigma\le\varsigma_{\max}$,
\[
\bar\mu:=\mathcal S_G\sinh(L\varsigma),
\]
where $\mathcal S_G$ is the function of Section~\ref{ssec:lift-legs} for the reference monomial $a$. On $\mathcal N_P$ only
the four monomials of $P$ among those free of $Z$ and $Z'$ have positive weight, so $\mathcal S_G=1+|\tilde\chi_1|+
|\tilde\chi_2|+|\mu_P\tilde\chi_1\tilde\chi_2|$ is a smooth function of $(\chi_1,\chi_2)$, at most $Ct^{-2\theta_2}$ within
$\theta_4$ of the node point; since $\varsigma\le\theta_2/4$ and $2\theta_2+\theta_2/4<1/2\le\theta_P$, the bound
$\bar\mu\le t^{-2\theta_P}/4$ of Definition~\ref{def:nodepiece} holds for $t$ small. With this choice the node piece and the
divisor pieces along the two legs of its passage are the same subset of $W^*$ where they overlap: by
Lemma~\ref{lem:NQ}(1), over $\mathscr C_P$ the node piece $\Pi^+_q$ is $\{Z'=0,\ (\chi_1,\chi_2)\in\mathscr C_P,\
|\hat Z|^2\le\mathcal S_G\sinh(L\varsigma)\}$, which is the divisor piece of depth $\varsigma$ of Lemma~\ref{lem:pieces}(1)
along either leg of the passage; that divisor piece is defined by the same inequality, which does not depend on the
reference monomial, and with the same ratio $\nu_Z/\nu_{Z'}$, which is constant within $2\theta_4$ of the node point.

We write $\Pi^\varsigma_q(g_{ab})$ and $\Pi^\varsigma_q(g_{bc})$ for $\Pi^+_q$, with this $\bar\mu$, over the branch of the
$b-a$ and of the $c-b$ passage, oriented as the lift of the unit germs $g_{ab}$ and $g_{bc}$ of (M1): by the orientation of
the germ followed by that of its torus (Section~\ref{ssec:nodepiece}).

\begin{lemma}[node pieces with multiplicity]\label{lem:nodesum}
For integers $c_{ab},c_{bc}$ the part of $\partial\bigl(c_{ab}\,\Pi^\varsigma_q(g_{ab})+c_{bc}\,\Pi^\varsigma_q(g_{bc})\bigr)$
over the thimble $\mathscr T_q$ is $(c_{ab}-c_{bc})\,\Sigma_q$, with $\Sigma_q$ oriented by Definition~\ref{def:orient}.
\end{lemma}

\begin{proof}
Both pieces contain $\Sigma_q=\pi^{-1}(\{0\}\times\mathscr T_q)$ (Lemma~\ref{lem:NQ}(3)). The thimble $\mathscr T_q$ is
symmetric in $x$ and $y$, so it is the same disc for the two passages, and $\Sigma_q$ and its orientation do not depend on
the passage (Definition~\ref{def:orient}). The unit germs are standard halves carrying the covectors
$-\langle\,\cdot\,,b-a\rangle$ and $+\langle\,\cdot\,,c-b\rangle$, with node coefficients $\rho_q(g_{ab})=1$ and
$\rho_q(g_{bc})=-1$ (Section~\ref{ssec:lift-model}); by Lemma~\ref{lem:SI} the part of $\partial\Pi^\varsigma_q(g)$ over
$\mathscr T_q$ is $\rho_q(g)\,\Sigma_q$ for each of them. The statement follows by adding.
\end{proof}

In the normal form the summands containing the node halves of $g_{ab}$ and $g_{bc}$ carry $c_{ab}$ and $c_{bc}$ times
the covectors of the unit germs, where $b^M_{m_P}=c_{ab}\,\partial g_{ab}+c_{bc}\,\partial g_{bc}$ in the basis of
Lemma~\ref{lem:unitgerms}; both passages may occur, with any integer coefficients, and then $c_{ab}-c_{bc}=\rho_q(b^M)$.
The two pieces overlap in $W^*$, over $\mathscr T_q$ and along the $S^1_\alpha$-orbits over the half-axis $C^M(u',P)$; as
chains they add, and no matching between them is required. Lemma~\ref{lem:SI} is applied only to the two unit germs, and
the multiplicities enter by adding. It covers both halves, both orientations and both signs of the covector for either
passage, so the other choices of side in (M1), which change the normal form by transverse discs, give the same boundary
over the thimble.

\subsection{Crossings with an arbitrary invariant coefficient}\label{ssec:lift-cross}

A transverse disc $v_A\mathfrak N_A$ of the normal form may carry any nonzero $v_A\in\operatorname{Inv}^*(A)$: a
non-primitive covector, a covector whose kernel meets the vanishing direction with any multiplicity, or a multiple of
$d^*_A$ itself. The following lemma gives its piece in all cases. We use the chart $U_{u,u'}$ with the coordinates
$Z=\chi^{w_u}$, $Z'=\chi^{w_{u'}}$, $\chi_1=\chi^{n'-n}$, $\chi_2=\chi^{n_y}$ of Section~\ref{ssec:lift-legs}, the dual basis
$(e_Z,e_{Z'},e_{\chi_1},e_{\chi_2})$ of $M$ to the basis $(w_u,w_{u'},n'-n,n_y)$ of $N$, and, for $\varsigma>0$, the normal
disc $\mathfrak D_\varsigma:=T_e^{-1}(\{y=y(\mathfrak x_A),\ y_e^2+\alpha_e^2\le\varsigma^2\})\subset B_Y$, the disc
$D_\varsigma:=E_t(\{y_e^2+\alpha_e^2\le\varsigma^2\})$ in the leg model base, and the Lefschetz neighbourhood
$U_{D_\varsigma}:=\{(Z,Z',\chi_1):1+\kappa_e\chi_1+\mathfrak bZZ'=0,\ (x_1,x_2)\in D_\varsigma\}$, with
$\log|\chi_2|=Ly(\mathfrak x_A)$ fixed.

\begin{lemma}[crossings]\label{lem:crossing}
Let $v\in\operatorname{Inv}^*(A)$, $v\ne0$. Write $v=mv_0$ with $m\ge1$ and $v_0$ primitive, and $v_0=k_g\,g^*_A+k_d\,d^*_A$.
Put $\Lambda_{v_0}:=\ker v_0\cap\Lambda^M$ and $\mathbb T(v_0)=\Lambda_{v_0}\otimes U(1)$.
\begin{enumerate}
\item $u'-u\in\Lambda_{v_0}$, and $\Lambda_{v_0}=\langle u'-u,\ell_0\rangle$ with $\ell_0:=k_g\hat e_1-k_d\hat e_2$, where $\hat
e_1:=a_1e_Z+e_{\chi_1}$, $\hat e_2:=a_2e_Z+e_{\chi_2}$ and $n=-w_u-w_{u'}+a_1(n'-n)+a_2n_y$. Put $k_1:=\langle
\ell_0,n'-n\rangle=k_g$ and $k_2:=\langle\ell_0,n_y\rangle=-k_d$; then $\gcd(k_1,k_2)=1$. The monodromy of $\Lambda^M$ around
$A$ preserves $\Lambda_{v_0}$ and acts on it, in the basis $(u'-u,\ell_0)$, by $\bigl(\begin{smallmatrix}1&k_g\\0&1
\end{smallmatrix}\bigr)$.
\item If $k_g\ne0$, the set
\[
\Pi(v_0):=\{(Z,Z',\chi_1,\chi_2):(Z,Z',\chi_1)\in U_{D_\varsigma},\ \log|\chi_2|=Ly(\mathfrak x_A),\ \chi_2^{k_g}\chi_1^{k_d}\in
\R_{>0}\}
\]
is a connected $|k_g|$-sheeted cover of $U_{D_\varsigma}$, a smooth four-manifold without boundary over $\mathfrak x_A$, with
fibre of type $I_{|k_g|}$ over $\mathfrak x_A$ and fibre $\mathbb T(v_0)$ over the other points of $\mathfrak D_\varsigma$; over
$\partial\mathfrak D_\varsigma$ it is the union of the $\mathbb T(v_0)$-orbits through the section of phase $0$.
\item If $k_g=0$, then $v_0=\pm d^*_A$ and $\mathbb T(v_0)=\mathbb T_d$, the exact torus of the leg. It is generated by the
rotation $u'-u$, which rotates $Z$ and $Z'$ oppositely, and by $\hat e_2$, which rotates $\chi_2$ and $Z$; it fixes
$\chi_1$, and it acts freely at every point with $Z,Z',\chi_2\ne0$. The set $\Pi(v_0):=\mathbb T_d\cdot\mathfrak s^\vee
(\mathfrak D_\varsigma)$ is homeomorphic to $\mathfrak D_\varsigma\times T^2$, and its boundary lies over
$\partial\mathfrak D_\varsigma$, where it is the union of the $\mathbb T_d$-orbits through $\mathfrak s^\vee$.
\end{enumerate}
The \emph{crossing piece} of $v\mathfrak N_A$ is $m\,\Pi(v_0)$, with $\varsigma=\varsigma_0$, oriented as the lift of the
disc: the orientation of $\mathfrak N_A$ followed by that of $\mathbb T(v_0)$.
\end{lemma}

The proof is in Appendix~\ref{app:liftproofs}.

For a strict tropical $2$-cycle of \cite{CBM}, the covector at a crossing is primitive and $|k_1|=1$
\cite[Def.~7.2(i)]{CBM}; then (2) is the piece of \cite[\S7.7, Step~6]{CBM}. For relative cycles no splitting of $v$ is
needed: when $k_g\ne0$ the piece $\Pi(v_0)$ is used directly, for every $|k_g|\ge1$, and the only additional case is $k_g=0$.

\subsection{Local pieces and collar paths}\label{ssec:lift-collar}

We truncate the local lifts and join their boundary tori to the positive real section by collar paths.
The estimates below keep these paths in the regions used to extend the lift over the regular part.

\subsubsection*{Truncation near the discriminant}
Depths are defined per summand of the normal form~\eqref{eq:nf2} near $\Gamma^M$. Fix for every $\Gamma^M$-vertex $p$
a closed ball $B^0_p$ of radius $r_B\le r_0$ about $p$, inside the neighbourhoods of Lemma~\ref{lem:place}(1)--(3) and
meeting $\Gamma^M$ only in the arcs at $p$. Inside $B^0_p$: a node half has depth $|\alpha_q|$; a tripod at a model vertex
$m$ with edge $[u,u']$ has depth $\alpha_{uu'}-\alpha_{uu'}(m)$; a turning germ has depth equal to the smaller of the
offsets from its two arcs, which are affine coordinates near the turning point, since the angle between the arcs there is
less than $\pi$ (Lemma~\ref{lem:arcsystems}). Outside the balls: a sheet along an arc (exit strip, reference sheet) has
depth $|\alpha_e|$, the offset from its arc in the chart $T_e$; a rotating stretch and a transverse disc have depth
$(y_e^2+\alpha_e^2)^{1/2}$. These descriptions agree where they overlap at the points of depth at most $2\varsigma_0$
(with $\varsigma_0$ as below): at the model nodes and vertices exactly, since the model arcs begin with level segments of
$\alpha_{uu'}$, where $\bar\alpha$ is constant; at the ends of the stretches since $\Theta=\pm\pi/2$ there; and near
$\partial B^0_p$ at a turning point because there a point at offset at most $2\varsigma_0$ from one arc is at offset at
least $c\,r_B-C\cdot2\varsigma_0$ from the other, where $c>0$ depends only on the angles between the arcs at the turning
points and $C\ge1$ is the largest ratio of the gradient norms of the two offsets there. Put $c_B:=c/(C+1)$; if
$\varsigma_0\le c_Br_B/8$, then $c\,r_B-2C\varsigma_0>2\varsigma_0$, so the smaller offset is the offset from the arc of the
strip.

\begin{definition}[depth]\label{def:depth}
For a summand $g$ and $0<\varsigma\le2\varsigma_0$, $g_{\le\varsigma}$ is the set of points of $g$ of depth at most
$\varsigma$, the \emph{inner boundary} $C^{(g)}_\varsigma$ is its frontier in $g$ away from the junction segments, and the
\emph{regular part} of $g$ is $g_{\ge\varsigma}:=\overline{g\setminus g_{\le\varsigma}}$. The inner boundary of a tripod has a
trivalent point on the half-edge $C^M(u',\omega)$.
\end{definition}

Fix a metric on $B_Y$ and a common Lipschitz constant $C_T\ge1$ of the inverse charts $T_e^{-1}$, $T_q^{-1}$ and of the
inverse depth parametrisations near the model vertices and turning points. Fix
\[
\varsigma_0\le\min\bigl(\varsigma_{\max}/2,\ r_S/4,\ r_\natural/(4C_T),\ c_Br_B/8\bigr).
\]
A point at depth at most $2\varsigma_0$ then lies within $2C_T\varsigma_0\le r_\natural/2$ of $\Gamma^M$; so every piece of
depth at most $2\varsigma_0$, and every collar strip below, lies over the region where $z^\natural$ consists exactly of the
germs and the normal discs (Lemma~\ref{lem:nfY}), inside the polygonal cylinders of the smoothed rotations.

\emph{Base maps.} The local pieces of depth $\varsigma$ (Lemmas~\ref{lem:pieces} and~\ref{lem:crossing}, and the node
pieces $\Pi^\varsigma_q$ of Section~\ref{ssec:lift-node}) lie over the model surfaces: for a summand $g$ the base map $\pi_g$ sends a point $x$
to the point of the model surface with leg coordinate $y(x):=\log|\chi_2(x)|/L$ (for the arcs of a node passage we take
$n_y=d-a$, respectively $n_y=b-a$, according to the passage of $g$, so that $y$ is the coordinate $y_2$, respectively $y_1$, of
Lemma~\ref{lem:nodemodel}(4) up to an additive constant), and with the tropical coordinates
$(y_e,\alpha_e)=E_t^{-1}(\log|\kappa_e\chi_1(x)|,\hat\mu(x))$ on $\Pi_{\mathfrak a,\varsigma}$ and on the crossing
pieces, and the offset $\frac1L\operatorname{arsinh}(\hat\mu(x)/\mathcal S_G(x))$ from the arc on the divisor pieces, the
node pieces and the vertex pieces. Near a model vertex, the position on the three sectors and the half-edge is given by a
fixed retraction, onto the tripod of the three arcs, of a neighbourhood of it in the plane of the torus coordinates of
$\omega$, which contains the image of the pair of pants; near a turning point, the offsets from the two arcs are given by
the moduli of the two Cox coordinates $z_2,z_3$. Where two of these descriptions overlap they differ by $O(1/L)$, and
$\pi_g$ is obtained by interpolating them with a partition of unity in the tropical coordinates; since
$E_t^{-1}(0,x_2)=(0,\frac1L\operatorname{arsinh}(x_2/2))$ and $\mathcal S_G=2$ on the branch over $\Omega_e$, the
descriptions agree to this order.

\emph{The piece of a summand.} For a summand $g$ and $0<\varsigma\le2\varsigma_0$ let $\Pi^\varsigma(g)$ be the set of
points $x$ of the local pieces of depth $\varsigma$ along $g$ with $\pi_g(x)\in g$: the node piece $\Pi^\varsigma_q$ over the
branch of its passage near a model node, the vertex piece $\Pi_v$ near a model vertex, the turning piece in $D_{u_i}$
near a turning point if $g$ contains $\mathrm{Turn}(u_i)$, the divisor pieces along its exit strips and reference sheets,
the pieces $\Pi_{\mathfrak a,\varsigma}$ along its rotating stretch, and the crossing piece if $g$ is a transverse disc. The
local pieces agree where they overlap (Lemma~\ref{lem:pieces}(1) and Section~\ref{ssec:lift-node}), so $\Pi^\varsigma(g)$ is
one subset of $W^*$, each point counted once, and $\pi_g$ maps it onto $g_{\le\varsigma}$, a homeomorphism up to the collapse
of the torus orbits. If the coefficient of $g$ is $m\,v_0$ with $m\in\Z$ and $v_0$ primitive (for a tripod see below), the \emph{piece} of $g$ is
$m\,[\Pi^\varsigma(g)]$, with $\Pi^\varsigma(g)$ oriented as the lift of $g$: the orientation of $g$ followed by that of
$\mathbb T(v_0)$. The piece of $-v_0$ is minus that of $v_0$, since the set is the same and the torus orientation is
reversed. The \emph{outer part} of the piece is $\pi_g^{-1}(C^{(g)}_\varsigma)\cap\Pi^\varsigma(g)$. For a summand containing
the node half of $g_{ab}$, respectively $g_{bc}$, the part of the piece over the node cube is $c_{ab}\,\Pi^\varsigma_q(g_{ab})$,
respectively $c_{bc}\,\Pi^\varsigma_q(g_{bc})$, restricted to base points in $g$; the thimble lies over $m_P\in g$. A tripod
summand of multiplicity $m$ carries $m\langle\,\cdot\,,e_k\rangle$ on its $k$-th sector; its piece is $m\,\Pi_v$ near the
vertex, continued by the divisor pieces of its exit strips, oriented sector-wise as the lift: over the $k$-th sector by the
orientation of that sector followed by that of $\mathbb T(\langle\,\cdot\,,e_k\rangle)$, and over the half-edge, where its
fibre is the filling $\mathcal Q=S^1_\alpha\times\mathcal Q'$ of Lemma~\ref{lem:pieces}(2), with the sign convention of
\cite[\S7.7, Step~2]{CBM}; $\Pi_v$ is counted once; at a turning point the pieces of the summands
containing $\mathrm{Turn}(u_1)$ and $\mathrm{Turn}(u_2)$, the latter oriented by $-o_e$ as in (M3), lie in the different
divisors $D_{u_1}$ and $D_{u_2}$ and are separate chains. For example, a turning point with away-oriented leg data
$(c_{12},c_{13},c_{23})=(1,1,-2)$, which arises from a positive vertex of $B'_\sigma$ with two turns ending on the same leg,
is $-\mathrm{Turn}(u_1)+2\,\mathrm{Turn}(u_2)$ in the basis $(\mathrm{Turn}(u_1),\mathrm{Turn}(u_2))$, or equally
$\mathrm{Turn}(u_2)+\mathrm{Turn}(u_3)$, the two differing by the full disc $\mathrm{Turn}(u_1)-\mathrm{Turn}(u_2)+
\mathrm{Turn}(u_3)$, which has zero leg data; either way the lift consists of separate turning pieces whose exit sheets
run along the common arc on opposite sides, one of them rotated to the reference side. Since the summands are
separate, their inner boundaries and pieces may cross in $B_Y$ and in $W^*$; nothing below refers to such crossings.

\subsubsection*{Collar paths}
Every piece of Lemma~\ref{lem:pieces}, and every node piece, has phase $\pi$ for the character $\chi^{\tilde v}$ of its
covector $v$, up to $O(\eta_0^{1/2})$ near the nodes, while the union of the $\mathbb T(v)$-orbits through $\mathfrak s^\vee$
has phase $0$. On the strip $g_{\le2\varsigma_0}\cap g_{\ge\varsigma_0}$ of a summand $g$ the lift consists of
\emph{collar paths}, paths of $\mathbb T(v)$-orbits in the exact models over the strip. Over the transverse segment from
$b_0\in C^{(g)}_{\varsigma_0}$ to $b_1\in C^{(g)}_{2\varsigma_0}$ the collar path first runs through the pieces of depths
between $\varsigma_0$ and $2\varsigma_0$ over that segment, at phase $\pi$; then, at fixed $\hat\mu$ and fixed moduli of the
torus coordinates, which determine the point up to the action of $S^1_\alpha$ once the phases of the torus coordinates are
given (Lemma~\ref{lem:quot} and its analogue for the leg model; where two monomials are in transition, their weights depend
on $|\hat Z\hat Z'|$ with logarithmic derivatives $O(1/L)$, and $|\hat Z\hat Z'|$ is the fixed point of a contraction), it
turns the phase of $\chi^{\tilde v}$ from $\pi$ to $0$ along one of the two arcs of the circle joining them, ending at a
point of $\Sigma^+_W$; finally it moves in $\Sigma^+_W$ to $\mathfrak s^\vee(b_1)$, along the image under $\mathfrak s^\vee$
of a segment of length $O(1/L)$. The two choices of the arc differ by the fibre loop $\gamma_v$ with $\langle v,\gamma_v\rangle
=1$. The collar path of a summand with coefficient $m\,v_0$ is $m$ times that of $v_0$. Around a crossing no phase is
turned, since the crossing piece has phase $0$: the collar path runs in $\Sigma^+_W$ from the positive point of the piece
over $b_0$ to $\mathfrak s^\vee(b_1)$, and the union of the torus orbits through these paths is the crossing piece over the
strip.

\emph{Direction convention.} Call two arcs of $\Gamma^M$ \emph{linked} if they are the two arcs of a turning germ, or the
two legs of one passage at a model node, and call the classes of the equivalence relation they generate \emph{linkage
classes}. All arcs of a class carry the same covector $d^*_A$ up to sign: through a turning point all arcs of $E_e$ have
$d^*_A=\pm\langle\,\cdot\,,n'-n\rangle$, and through a node passage $b-a=c-d$. So a class determines a set of edges of
$\mathcal F_C$, closed under the translations $e\mapsto e+(d-a)$ and $e\mapsto e+(b-a)$ across node parallelograms, and a
translation preserves the vector of an edge. Orient these edges compatibly and let $d_{\rm cl}:=n'-n$ for the oriented
edge $[n,n']$; no cycle of translations reverses the orientation. At an $A_2$ hexagon the two legs of a node that lie on
node-to-node arcs are adjacent and belong to different passages, so a linkage class enters the hexagon along an outer
leg, passes two of its nodes and leaves along an outer leg; it does not run around the cycle of node-to-node arcs. On
every arc of a class, and for every summand other than a transverse disc, the geometric collar path turns the phase of $\chi^{d_{\rm cl}}$ from $\pi$ to $0$ in the positive
direction; for a summand with covector $m(\pm d^*_A)$ the collar is $\pm m$ times this geometric family, which does not
depend on the sign, since the phases $\pi$ and $0$ of $\chi^{\pm d_{\rm cl}}$ coincide.

\begin{lemma}[collar paths]\label{lem:collarB}
\begin{enumerate}
\item Along every summand, and through a turning germ or a node passage, the direction of the collar paths is
constant.
\item At a model vertex the three directions are independent. For a tripod of multiplicity $m$ the collar cells
$\mathcal Q\times(\text{transverse interval})$ below are taken $m$ times.
\item At every junction segment the collars of the adjacent summands cancel.
\item Around a crossing no phase is turned, and the collar runs in the positive locus.
\item A different choice of the direction on a linkage class changes the chain by a closed chain.
\end{enumerate}
\end{lemma}

\begin{proof}
(1) The strip cells along an arc are rectangles (segment $\times$ collar path) filled by the explicit family; a change of
direction would leave the loop class $\pm\gamma_v\ne0$ on the boundary of a rectangle. Through a turning point and through a
node passage the character is the same on both arcs, by the direction convention. The two passage legs are joined inside
the branch $\mathscr C_P$, along which the phase of $\tilde\chi_1$ is $\pi$ up to $O(r)$, so the rectangle over the node
contributes no winding: the collar paths through the node form a homotopy of cylinders from $\mathscr C$ to
$\{\arg\tilde\chi_1=0\}$, first $\mathscr C\simeq\{\tilde\chi_1=-\mu_P^{-1}\}$, since the two capping discs (the thimble
with the inner half-annulus of $\mathscr C_P$, and $\{x=0,\ |y|\le r\}$) have the same boundary up to a small isotopy and
lie in the Morse bidisc, hence are homotopic rel boundary in the region of Lemma~\ref{lem:R1}; then the phase of
$\tilde\chi_1$ is rotated from $\pi$ to $0$. At $\hat\mu>0$ the $S^1_\alpha$-orbits are free, so the homotopy lifts. The
second circle of the node, the rotation of $\tilde\chi_2$ for the $b-a$ passage or of $\tilde\chi_1$ for the $c-b$
passage, lies in the fibre torus of the germ, so no half-turn correction arises.
(2) The only cells joining the three strips at a model vertex are the fillings $\mathcal Q$ over the half-edge times the
transverse interval of the strip. A filling of three circles with $\sum_ks_k[G_k]=0$ in $T^3/S^1_\alpha$ exists for every
triple of phases, including the concurrent one at $\mathfrak s^\vee$, so every path of triples lifts.
(3) The collars on the two sides of a junction segment are given by the same formula, with the same direction and
covectors that are multiples of $d^*_A$ summing to zero.
(4) The crossing pieces have phase $0$ (Lemma~\ref{lem:crossing}), so the collar paths there run in $\Sigma^+_W$, as
described above.
(5) Changing the direction adds the union of the torus orbits through a full phase turn over the strip, a closed chain.
\end{proof}

By Lemma~\ref{lem:interface} below, the collar paths of a summand over a cell of its strip lie in the region
$W^*_{V_\kappa}$ associated with the adjacent simplex $\kappa$ of the triangulation $\mathfrak T$
(Section~\ref{ssec:lift-regular}); these regions are defined in Definition~\ref{def:regionsB}.

\subsection{Extending the lift over the regular part}\label{ssec:lift-regular}

The local pieces determine boundary data for a cell-by-cell extension over the remaining domain.
The essential input is Lemma~\ref{lem:R1}: each target region is aspherical, and its fundamental group is the prescribed integral lattice.

\subsubsection*{Regions}
Over the regular part the lift is built in regions of $W^*$ defined by shadows. For a vertex $u$ of $\mathcal T$ let
$\operatorname{Vis}(u)\subseteq B_Y$ be the union of the facets $F^*_n$ with $n\in u^\vee$, and $\pi_u\colon M_\R\to
M_\R/\R u$ the projection. Then $\pi_u$ is injective on $\operatorname{Vis}(u)$: if $p\in F^*_n$ and $p+su\in F^*_{n'}$ with
$n,n'\in u^\vee$, then $0\ge\operatorname{val}_n(p+su)=-s$ and $0\ge\operatorname{val}_{n'}(p)=s$. The vertex region $R^M_u$
lies in $\operatorname{Vis}(u)$. Let $U_u\subset v_{\mathcal T}$ be the union of the dense torus and the open orbit of $D_u$.
The characters $\chi^m$, $m\in u^\perp\cap N$, are regular and nowhere zero on $U_u$, so $\pi_u(\operatorname{Log}x):=
(\log|\chi^m(x)|)_m$ is defined on $U_u$. For $x\in W^*\cap U_u$ with $\pi_u(\operatorname{Log}x)/L\in\pi_u(\operatorname{Vis}
(u))$, the \emph{$u$-shadow} $\operatorname{sh}_u(x)$ is the point of $\operatorname{Vis}(u)$ with $\pi_u(\operatorname{sh}_u
(x))=\pi_u(\operatorname{Log}x)/L$; for $x$ in the torus, $\operatorname{Log}x/L=\operatorname{sh}_u(x)+\lambda_u(x)\,u$.
For a vertex $n$ of $\mathcal F_C$ fix $m_n\in M$ with $\langle m_n,n\rangle=1$; the projection $\varpi_n\colon M_\R\to M_\R/
\R m_n$ restricts to an affine isomorphism of the hyperplane $\{\operatorname{val}_n=0\}$. For $x$ in the torus the
\emph{$n$-shadow} $\operatorname{sh}_n(x)$ is the point of that hyperplane with the same image under $\varpi_n$ as
$\operatorname{Log}x/L$, and $\operatorname{Log}x/L=\operatorname{sh}_n(x)+\lambda_n(x)\,m_n$.

\begin{definition}[regions]\label{def:regionsB}
A \emph{regular set} is an open set $V\subseteq B_Y$ such that either its closure is compact in the open star of a cell
$C^M(u,\omega)$ and $\pi_u(V)$ is convex ($V$ is of \emph{type $u$}), or $V$ is contained in no vertex region, its closure is
compact in the open star of a cell $C^M(\tau,n)$ with $\dim\tau\ge1$, and $\varpi_n(V)$ is convex (type $n$). For
$\epsilon>0$ the \emph{region} over $V$ is
\begin{align*}
W^*_V&:=\{x\in W^*\cap U_u:\ \operatorname{sh}_u(x)\in V,\ \lambda_u(x)<\epsilon\}&&(\text{type }u),\\
W^*_V&:=\{x\in W^*\cap T:\ \operatorname{sh}_n(x)\in V,\ |\lambda_n(x)|<\epsilon\}&&(\text{type }n),
\end{align*}
where $\lambda_u:=-\infty$ on $D_u$ and $T$ is the dense torus.
\end{definition}

The open stars of the cells $C^M(u,\omega)$ and $C^M(\tau,n)$ cover $B_Y\setminus\Gamma^M$ (Proposition~\ref{prop:L1}); the
first lie in $R^M_u$, the second in $\operatorname{int}F^*_n$.

\begin{lemma}[regular regions]\label{lem:R1}
Let $V$ be a regular set. There is $\epsilon_V>0$ such that for $\epsilon\le\epsilon_V$ and $t<t_V(\epsilon)$:
\begin{enumerate}
\item if $V$ is of type $u$, then $W^*_V$ is the graph of a smooth function $Z_u=\varphi$ over the set
$\operatorname{Log}^{-1}(L\,\pi_u(V))$ of the torus $(\C^*)^3$ of the other Cox coordinates of a maximal cone containing $u$;
so $W^*_V\cong T^3\times\pi_u(V)$, and the inclusion $W^*_V\subset U_u$ induces an isomorphism $\pi_1(W^*_V)\to H_1(U_u)=
M/\Z u$;
\item if $V$ is of type $n$, then $W^*_V$ is the graph of a smooth function $\chi^n=\varphi$ over the set
$\operatorname{Log}^{-1}(L\,\varpi_n(V))$ of the torus of the characters in $m_n^\perp\cap N$, and the inclusion $W^*_V\subset
T$ induces an isomorphism of $\pi_1(W^*_V)$ onto $n^\perp\cap M\subset M=H_1(T)$;
\item if $V'$ and $V$ are regular sets with $\overline{V'}\subset V$, and $W^*_{V'}$, $W^*_V$ are taken with constants
$\epsilon'\le\epsilon$, then $W^*_{V'}\subseteq W^*_V$ for $\epsilon\le\epsilon_{V',V}$ and $t<t_{V',V}(\epsilon')$, also when $V'$
is of type $u$ and $V$ of type $n$; the inclusion induces the parallel transport of $\Lambda^M$.
\end{enumerate}
In particular every region is aspherical, with fundamental group the lattice of flat sections of $\Lambda^M$ over $V$.
\end{lemma}

The proof is in Appendix~\ref{app:liftproofs}.

\begin{lemma}[inner boundary]\label{lem:interface}
Let $\varsigma:=\varsigma_0$, and let $x$ be a point of $W^*$ of one of the following kinds: a point of the piece
of depth $2\varsigma$ of a summand $g$ whose base point $\pi_g(x)$ lies at depth at least $\varsigma$; a point of a
collar path of $g$ over a point $b$ of its strip, and then $\pi_g(x):=b$; or a point $\mathfrak s^\vee(b)$, and then
$\pi_g(x):=b$. There are $\varsigma_1\in(0,\varsigma/2]$ and $C$, depending only on the configuration, the model arc
systems, $\Upsilon^M$ and $\varsigma$, such that:
\begin{enumerate}
\item if $L\ge1/\varsigma_1$, then for every vertex $u$ of $\mathcal T$ with $\pi_g(x)\in R^M_u$ the point $x$ lies in $U_u$,
$|\operatorname{sh}_u(x)-\pi_g(x)|\le C/L$, and $\lambda_u(x)\le C/L$ if $x$ lies in the torus;
\item for every $\varsigma_2>0$ there is $L_{\varsigma_2}$ such that, if $\pi_g(x)$ lies at distance at least $\varsigma_2$ from
the two-skeleton of $B_Y$ and $L\ge L_{\varsigma_2}$, then $x$ lies in the torus and $|\operatorname{Log}x/L-\pi_g(x)|\le C/L$.
\end{enumerate}
Consequently, if $V$ is a regular set and $\mathcal W\subset V$ is compact, then every such $x$ with $\pi_g(x)\in\mathcal W$
lies in $W^*_V$, for $t<t_{V,\mathcal W}$.
\end{lemma}

The proof is in Appendix~\ref{app:liftproofs}.

\subsubsection*{The abstract domain}
Let $\mathscr S$ be the finite $2$-complex obtained from the disjoint union of
\begin{itemize}
\item every summand of~\eqref{eq:nf2} near $\Gamma^M$ (the transverse discs $\mathfrak N^{\varepsilon'_A}_A$ included),
restricted to depth at least $2\varsigma_0$ and continued inside $U$ up to $\partial U$,
\item the correction strips and annuli of the two local replacements,
\item the simplices of $y^\natural$ outside $U$,
\end{itemize}
by identifying (i) the junction segments of adjacent summands, beyond depth $2\varsigma_0$, and (ii) the curves along
which a summand ends on another summand: on $\partial U$, where it meets $y$; on the circles
$\partial\mathfrak N^{\varepsilon'_A}_A$ and $\partial(\bar D(E_A)\cap N_{\varepsilon_A})$ bounding the annuli; and on the
polygonal cylinders $\{\varrho_{\mathsf P}=\varrho_0\}$, where it meets a correction strip (this includes the seam of a
smoothed rotation, within the stretch after $y^*_{A,p}$). Segments interior to a summand, such as the transverse
segments where a germ continues as its own exit strip and the start of a rotation, are not cut. Summands that meet
transversally in $B_Y$ (a transverse disc and a sheet along its ray, the two node halves along $C^M(u',P)$, exit sheets of
different germs meeting only on $\Gamma^M$) are \emph{not} identified. The triangulation $\mathfrak K_2$ is chosen so that
the finitely many curves of transversal meeting cross $\partial U$ and the boundaries of the replacement neighbourhoods and
cylinders transversally in finitely many points. Triangulate $\mathscr S$ compatibly, each summand triangulated so that
the triangulations agree on the identified curves and $y^\natural$ simplicially, with the inner boundary curves
$C^{(g)}_{2\varsigma_0}$ as subcomplexes, whose vertices are the ends of the transverse collar segments issuing from the
vertices of a fixed triangulation of $C^{(g)}_{\varsigma_0}$ (the strips $g_{\le2\varsigma_0}\cap g_{\ge\varsigma_0}$ themselves are
not part of $\mathscr S$; they are lifted by the collar paths). The
inclusions define a continuous map $j^\natural\colon\mathscr S\to B_Y\setminus\Gamma^M$. Put on each $2$-simplex $\kappa$ of
$\mathscr S$ the coefficient $v_\kappa$ of its summand, a flat section of $(j^\natural)^*(\Lambda^M)^\vee$ over $\kappa$.

\begin{lemma}[the coefficient chain]\label{lem:Xichain}
$v=\sum_\kappa v_\kappa[\kappa]$ is a $2$-chain on $\mathscr S$ whose boundary is supported on the inner boundary curves, where it is
$\sum_g(\text{coefficient of }g)\,[C^{(g)}_{2\varsigma_0}]$.
\end{lemma}

\begin{proof}
At an edge inside a summand, including the transverse segments where a germ continues as its own exit strip and
the start of a rotation, the two sides carry the same coefficient. At a tripod edge, on the half-edge of a tripod
summand of multiplicity $m$, the three adjacent coefficients $m\langle\,\cdot\,,e_k\rangle$ sum to zero with the
orientations of (M2). At a junction segment the adjacent coefficients
are multiples of one covector $d^*_A$ summing to zero. At an identified curve of type (ii) the adjacent triangles of
$\mathscr S$ are exactly the triangles of $z^\natural$ at that edge of $B_Y$, by the transversality of the meeting curves,
and their coefficients sum to zero because $z^\natural$ is a relative cycle with boundary only on $\Gamma^M$.
\end{proof}

Since summands that meet transversally are different components of $\mathscr S$ near the meeting curve, $\mathscr S$
has no edges along which two summands cross: a transverse disc and a sheet on the same arc, the two node halves at a
node used by both passages, and the two turning germs of a turning point used twice on one arc are lifted separately.
Branching edges inside $y^\natural$, coming from the arbitrary chain off $\Gamma^M$, remain; they are treated like any
other edge.

\subsubsection*{Triangulation and regular sets}
Let $\mathfrak T$ be a subdivision of the triangulation of $\mathscr S$, fine enough that the image of every simplex has
small diameter. By a Lebesgue number argument each $j^\natural(\kappa)$ lies at positive distance from the frontier of the open
star of a cell $C^M(u,\omega)$ or $C^M(\tau,n)$; if it lies in a vertex region $R^M_u$ we take the star of a cell
$C^M(u,\omega)$, and otherwise that of a cell $C^M(\tau,n)$. Let $V_\kappa$ be the preimage, under $\pi_u$ or $\varpi_n$, of the
convex hull of the $\mathfrak r_{\dim \kappa}$-neighbourhood of the image of $\kappa$, with radii $\mathfrak r_0<\mathfrak r_1<\mathfrak
r_2$, each small compared with the next and $\mathfrak r_2$ small compared with the mesh; then $V_\kappa$ is a regular set,
$V_\kappa\supset j^\natural(\kappa)$, and $\overline{V_{\kappa'}}\subset V_\kappa$ for $\kappa'\subsetneq \kappa$. We take $\epsilon$ and $t$ so small that
Lemmas~\ref{lem:R1} and~\ref{lem:interface} hold for these finitely many sets, with $\mathcal W$ the images of the
simplices. This uses only that $j^\natural$ is continuous and $|\mathfrak T|$ compact, not that $\mathfrak T$ triangulates an
embedded surface.

\subsubsection*{The domain over the regular part}
For a simplex $\kappa$ of $\mathfrak T$ let $\Lambda_\kappa$ be the lattice of flat sections of $\Lambda^M$ over $V_\kappa$, which is
$\pi_1(W^*_{V_\kappa})$ (Lemma~\ref{lem:R1}); for $\kappa'\subset \kappa$ the inclusion $W^*_{V_{\kappa'}}\subseteq W^*_{V_\kappa}$ induces the
transport $\Lambda_{\kappa'}\cong\Lambda_\kappa$. We orient $\Lambda_\kappa\otimes\R$ by the orientation of $B_Y$, and put
$\mathbb T^3_\kappa:=\Lambda_\kappa\otimes\R/\Z$, a $K(\Lambda_\kappa,1)$. For a primitive covector $v^0$ on $\Lambda_\kappa$ the
\emph{canonical torus} is $\mathbb T(v^0):=(\ker v^0\otimes\R)/\ker v^0\subset\mathbb T^3_\kappa$, with the tautological
identification $\pi_1(\mathbb T(v^0))=\ker v^0\subseteq\Lambda_\kappa$, oriented so that $\operatorname{or}(\ker v^0)\wedge n=
\operatorname{or}(\Lambda_\kappa\otimes\R)$ for $v^0(n)>0$. Then $[\mathbb T(v^0)]\mapsto v^0$ is the restriction of the
linear isomorphism $H_2(\mathbb T^3_\kappa)=\Lambda^2\Lambda_\kappa\to\Lambda_\kappa^\vee$, $a\wedge b\mapsto\det(a,b,\cdot)$: for an oriented
basis $(a,b)$ of $\ker v^0$, $\det(a,b,\cdot)$ vanishes on $\ker v^0$, is positive where $v^0$ is, and is primitive.

\emph{Three kinds of edges.} An edge $e$ of $\mathfrak T$ is a \emph{sheet edge} if the coefficients of all its adjacent
$2$-simplices are, after transport to $\Lambda_e$, integer multiples $\mu_\kappa\,v^0_e$ of one primitive covector $v^0_e$; by
Lemma~\ref{lem:Xichain}, $\sum_\kappa[\kappa:e]\mu_\kappa=0$ unless $e$ lies on an inner boundary curve. Edges inside a summand and on junction
segments are sheet edges. An edge is a \emph{tripod edge} if it lies on the half-edge $C^M(u',\omega)$ of a tripod
summand. All other edges, the branching edges of $y^\natural$ and the edges on identified curves of type (ii) where
several summands meet, are \emph{general edges}. We take $\mathfrak T$ so fine that every edge meeting an inner boundary curve is a
sheet edge or a tripod edge; the identified curves of type (ii) lie at depth at least $r_\natural\ge4\varsigma_0$.

\emph{Fibre complexes.} For a $2$-simplex $\kappa$ with $v_\kappa\ne0$ write $v_\kappa=m_\kappa\,v^0_\kappa$ with $m_\kappa\ge1$ and $v^0_\kappa$
primitive, and put $F_\kappa:=\mathbb T_\kappa:=\mathbb T(v^0_\kappa)$; if $v_\kappa=0$, $F_\kappa$ is a point. For $\kappa'\subset \kappa$ the transport maps
kernels to kernels and induces an isomorphism of canonical tori; these transports are induced by inclusions of regions, so
they compose consistently around every vertex of $\mathfrak T$, and no composite of identifications below is a nontrivial
automorphism of a torus. For an edge $e$:
\begin{itemize}
\item at a sheet edge, all adjacent $\mathbb T_\kappa$ are identified with $F_e:=\mathbb T(v^0_e)$ by the canonical maps, with the
orientation sign of $\mu_\kappa$, so that $m_\kappa[\mathbb T_\kappa]=\mu_\kappa[F_e]$; and $\mathscr Q_e:=0$;
\item at a tripod edge, of a tripod summand of multiplicity $m_e$, the three adjacent sectors carry the covectors
$m_e\langle\,\cdot\,,e_k\rangle$, whose kernels meet in a rank-one sublattice $\ell$ (the three $e_k$ span a plane). Let
$F_e:=\mathbb T^3_e$, with a cell structure in which the three subtori $\mathbb T_\kappa$ are subcomplexes, and
$\mathscr Q_e:=m_e\,\mathscr Q_e^0$, $\mathscr Q_e^0:=(\ell\otimes U(1))\times\mathcal Q'_e$, where $\mathcal Q'_e$ is a $2$-chain in
$\mathbb T^3_e/(\ell\otimes U(1))$ with $\partial\mathcal Q'_e=\sum_ks_kG_k$, the images of the three tori with the signs of the
boundary orientations \cite[\S7.7, Step~2]{CBM}. At the inner boundary choose this filling to be exactly the terminal
filling of the collar family in Lemma~\ref{lem:collarB}(2), and transport it along the subdivided tripod half-edge,
with the incidence orientation on each edge. This is a compatible choice: the branching half-edge of each tripod
summand meets its inner boundary once, and its other end lies on $\partial U$ beyond the collar region;
\item at a general edge, let $Y_e:=\bigvee_{\kappa\supset e}F_\kappa$, realise the labels by a cellular map $\gamma_e\colon Y_e\to\mathbb T^3_e$,
and let $F_e$ be its mapping cylinder, a cell complex containing $Y_e$ and retracting onto $\mathbb T^3_e$, labelled by
$\operatorname{lab}_e\colon\pi_1(F_e)\cong\Lambda_e$.
\end{itemize}
For a vertex $p$ of $\mathfrak T$, $F_p$ is the union of the $F_e$, $e\ni p$, glued along the common $F_\kappa$; the labels agree on
the $F_\kappa$ and define $\operatorname{lab}_p\colon\pi_1(F_p)\to\Lambda_p$, which kills the boundary of every $2$-cell, since each $2$-cell
lies in some $F_e$.

\emph{Fillings over general edges.} At a general edge, $\sum_{\kappa\supset e}[\kappa:e]\,v_\kappa=0$ in $\Lambda_e^\vee$ by
Lemma~\ref{lem:Xichain}. Under the linear isomorphism $H_2(F_e)=H_2(\mathbb T^3_e)\cong\Lambda^\vee_e$ above, the class of
$\sum_{\kappa\supset e}[\kappa:e]\,m_\kappa[\mathbb T_\kappa]$ is $\sum_\kappa[\kappa:e]\,m_\kappa v^0_\kappa=\sum_\kappa[\kappa:e]\,v_\kappa=0$, so there is a cellular $3$-chain
$\mathscr Q_e$ in $F_e$ with
\[
\partial \mathscr Q_e=\sum_{\kappa\supset e}[\kappa:e]\,m_\kappa\,[\mathbb T_\kappa].
\]
No common rank-one sublattice of the kernels and no bound on the valency of $e$ is needed.

\emph{The domain.} Let $\mathcal D^\vee_{\rm reg}$ be the union of the $|\kappa|\times F_\kappa$ over the simplices $\kappa$ of $\mathfrak T$,
glued along $|\kappa'|\times F_\kappa\subseteq|\kappa'|\times F_{\kappa'}$ for $\kappa'\subset \kappa$, with the fundamental chain
\[
[\mathcal D^\vee_{\rm reg}]:=\sum_{\kappa\in\mathfrak T_2}m_\kappa\,[\kappa\times\mathbb T_\kappa]+\sum_{e\in\mathfrak T_1\setminus\text{inner boundary}}[e\times \mathscr Q_e],
\]
each $\kappa\times\mathbb T_\kappa$ oriented as the product. Over an inner boundary edge $e$ the adjacent triangles lie in one summand
$g$, the other side of the inner boundary being the strip lifted by the collar paths, and the cells $e\times\mathbb T_\kappa$ are glued
to the ends of the collar paths over $e$, the union of the $\mathbb T_\kappa$-orbits through $\mathfrak s^\vee$; over the point
where the inner boundary meets a tripod edge, the filling $m_e\mathscr Q_e^0$ is glued to $m_e$ times the collar cells
$\mathcal Q\times(\text{transverse interval})$ of Lemma~\ref{lem:collarB}(2).

\begin{lemma}[boundary of the domain]\label{lem:Mbd}
Let $\check\Xi_{\rm loc}$ be the sum of the pieces of depth $\varsigma_0$ of the summands, with their multiplicities,
and of their collar paths (Section~\ref{ssec:lift-collar}). Then
$\partial(\check\Xi_{\rm loc}+[\mathcal D^\vee_{\rm reg}])$, with $[\mathcal D^\vee_{\rm reg}]$ glued as above, is
\[
\sum_q\rho_q(z)\,\Sigma_q+\sum_p\,[p\times z_p],\qquad
z_p:=\begin{cases}
\displaystyle\sum_{e\ni p}[e:p]\,\mathscr Q_e,&p\text{ off the inner boundary},\\
0,&p\text{ on the inner boundary},
\end{cases}
\]
the second sum over the vertices $p$ of $\mathfrak T$, where $z_p$ is a cellular $3$-cycle of $F_p$.
\end{lemma}

\begin{proof}
The piece of a summand $g$ is $m\,[\Pi^{\varsigma_0}(g)]$, a four-manifold with corners fibred over $g_{\le\varsigma_0}$ up
to the collapse of the torus orbits. The orbits collapse to circles over the arcs, and to points over the turning points, without
creating boundary (Lemma~\ref{lem:pieces}(1),(3)), and the crossing pieces have no boundary over $\mathfrak x_A$
(Lemma~\ref{lem:crossing}). The seams inside $g$, where the description near a vertex passes to the description along an
exit strip or where a rotation starts, contribute nothing: $\Pi^{\varsigma_0}(g)$ is one set with one coefficient, the local
descriptions being glued along their overlaps and not added; in particular a vertex piece $\Pi_v$ has, inside the piece
of its summand, only its outer part, and continues along its tentacles as the divisor pieces of the same summand.
So the boundary of the piece consists of three parts.
\begin{enumerate}
\item[(a)] Over the thimbles: for the two summands containing the node halves at $m_P$, the parts over the node cube are
$c_{ab}\Pi^{\varsigma_0}_q(g_{ab})$ and $c_{bc}\Pi^{\varsigma_0}_q(g_{bc})$, and by Lemma~\ref{lem:nodesum} their boundary over
$\mathscr T_q$ is $(c_{ab}-c_{bc})\Sigma_q=\rho_q(b^M)\Sigma_q=\rho_q(z)\Sigma_q$.
\item[(b)] Over the junction segments at depth at most $\varsigma_0$, which lie inside the polygonal cylinders at the far ends
of the rotations: the adjacent summands have pieces over the same subset of $W^*$ (divisor
pieces or pieces $\Pi_{\mathfrak a,\varsigma_0}$ of one arc, which agree across the junction), with coefficients multiples of
$d^*_A$ summing to zero with the orientation signs; these faces cancel.
\item[(c)] The outer parts over $C^{(g)}_{\varsigma_0}$. They are matched by the collar cells over the strips
$g_{\le2\varsigma_0}\cap g_{\ge\varsigma_0}$, which cancel along the junction segments (Lemma~\ref{lem:collarB}(3)) and end at
the unions of the torus orbits through $\mathfrak s^\vee$ over the curves $C^{(g)}_{2\varsigma_0}$; these are the cells
$e\times\mathbb T_\kappa$ over the inner boundary edges.
\end{enumerate}
Pieces of different summands that intersect in $W^*$ (the two node halves, a transverse disc and a sheet) are separate
chains, and their intersection contributes no boundary.

Over the regular part, $\partial(\kappa\times\mathbb T_\kappa)=\partial \kappa\times\mathbb T_\kappa$, so the faces over an edge $e$ are
$e\times\sum_{\kappa\supset e}[\kappa:e]m_\kappa\mathbb T_\kappa$. At a sheet edge off the inner boundary this is
$e\times(\sum_\kappa[\kappa:e]\mu_\kappa)F_e=0$. At a tripod edge and
at a general edge, $\partial(e\times \mathscr Q_e)=\partial e\times \mathscr Q_e-e\times\partial \mathscr Q_e$, and $e\times\partial \mathscr Q_e$ cancels the faces
over $e$ by the choice of $\mathscr Q_e$ (for tripod edges this is \cite[\S7.7, Step~2]{CBM}, multiplied by $m_e$). What remains
of the edge terms at a vertex $p$ off the inner boundary is $p\times z_p$, and
\[
\partial z_p=\sum_{e\ni p}[e:p]\sum_{\kappa\supset e}[\kappa:e]\,m_\kappa[\mathbb T_\kappa]=\sum_{\kappa\ni p}m_\kappa[\mathbb T_\kappa]\sum_{e}[e:p][\kappa:e]=0,
\]
because every $2$-simplex $\kappa\ni p$ contributes through its two edges at $p$ with opposite signs.
At an inner boundary vertex the boundary of the collar must also be included. Along a sheet it cancels the product
faces exactly. Where a tripod half-edge meets the inner boundary, its endpoint term $[e:p]\mathscr Q_e$ is cancelled by the
oppositely oriented terminal filling of the collar family, since these were chosen to be the same cellular chain.
There are no general edges at the inner boundary. Thus no residual chain remains there, giving $z_p=0$ as stated.
\end{proof}

\subsubsection*{The extension}
\begin{lemma}[extension]\label{lem:extB}
There is a continuous map $F^\vee\colon\mathcal D^\vee_{\rm reg}\to W^*$ which maps the cells over each simplex $\kappa$ of
$\mathfrak T$ into $W^*_{V_\kappa}$, sends fibre loops to loops representing their labels, is $\mathfrak s^\vee\circ j^\natural$ on the
base cells $\kappa\times\{\star\}$, and on the inner boundary cells agrees with the collar ends.
\end{lemma}

\begin{proof}
We extend cell by cell in order of dimension. On $0$-cells use $\mathfrak s^\vee$, which lies in $W^*_{V_\kappa}$ by
Lemma~\ref{lem:interface}. On $1$-cells, base cells go along $\mathfrak s^\vee$ and fibre $1$-cells to loops of the labelled
classes. A $2$-cell extends if its boundary loop is trivial in the abelian group $\Lambda_\kappa=\pi_1(W^*_{V_\kappa})$: for $\kappa\times\{\star\}$,
$\kappa$ a $2$-simplex, the boundary loop lies in $\mathfrak s^\vee(j^\natural \kappa)$, which is contractible in $W^*_{V_\kappa}$; for
$e\times s^1$, the two fibre loops represent the same parallel class (Lemma~\ref{lem:R1}(3)); for $p\times s^2$,
$s^2$ a $2$-cell of $F_p$, $\operatorname{lab}_p(\partial s^2)=0$. Cells of dimension three and four extend because $W^*_{V_\kappa}$ is
aspherical (Lemma~\ref{lem:R1}). The inner boundary cells are prescribed by the collars, which satisfy the same loop conditions
(Lemma~\ref{lem:collarB}) and lie in the right regions (Lemma~\ref{lem:interface}).
\end{proof}

\subsection{The positive real locus and the fibre class}\label{ssec:lift-fibre}
The extension leaves a sum of residual three-cycles, each homologous to a multiple of a torus fibre.
Intersection with the positive real three-sphere shows that their total fibre coefficient is zero.

Put $\check\Xi_0:=\check\Xi_{\rm loc}+F^\vee_*[\mathcal D^\vee_{\rm reg}]$, an integral four-chain on $W^*$. By
Lemma~\ref{lem:Mbd},
\[
\partial\check\Xi_0=\sum_q\rho_q(z)\,\Sigma_q+\sum_pF^\vee_*z_p .
\]
Each $F^\vee_*z_p$ lies in an aspherical region $W^*_{V_p}\simeq T^3$, so $[F^\vee_*z_p]=n_p[T^3_p]$ in
$H_3(W^*_{V_p})\cong\Z$, where $T^3_p$ is a fibre torus over $p$ oriented by the orientation of $\Lambda^M$. The monodromy of
$\Lambda^M$ lies in $SL(\Lambda^M)$, since that of $\Lambda_Y$ does (Proposition~\ref{prop:L2}(4)), and overlapping regions are
homotopy equivalent to each of their members; so all the $[T^3_p]$ agree in $H_3(W^*)$. We call this class $[T^3]$. It is
represented by an orbit of the compact subtorus $(n^\perp\cap M)\otimes U(1)$ through a positive point of a region over a
cell $C^M(\tau,n)$ at which only the monomials $0$ and $n$ have positive weight.

\begin{proposition}[no fibre-class boundary]\label{lem:F}
$\sum_pn_p=0$. Consequently there is an integral four-chain $\check\Xi^W(z)$ on $W^*$ with $\partial\check\Xi^W(z)=\sum_q
\rho_q(z)\,\Sigma_q$.
\end{proposition}

\begin{proof}
Replacing each $F^\vee_*z_p$ by $n_pT^3_p$ changes $\check\Xi_0$ by chains inside the regions $W^*_{V_p}$. Then
$\sum_q\rho_q\Sigma_q+\sum_pn_pT^3_p$ is a boundary, so $\sum_q\rho_q[\Sigma_q]+\bar n[T^3]=0$ in $H_3(W^*;\Z)$ with
$\bar n=\sum_pn_p$. We pair with the class of the positive locus $\Sigma^+_W$, a closed oriented three-manifold
(Lemma~\ref{lem:possec}).
\begin{itemize}
\item In the part of a region over a cell $C^M(\tau,n)$ where only $0$ and $n$ have positive weight,
$W^*=\{\gamma_0x^0+\gamma_nx^n=0\}$ is a translate of the subtorus $\ker\chi^n$, invariant under $\mathbb T_n:=(n^\perp\cap M)
\otimes U(1)$, and $\Sigma^+_W$ is its positive real part. At a positive point $x_b$ the orbit $\mathbb T_n\cdot x_b$ has
tangent space $i\R^3$ and $\Sigma^+_W$ has $\R^3$, complementary in $T_{x_b}W^*\cong\C^3$. An element of $\mathbb T_n$ that maps
$x_b$ to a positive point is the identity, so the intersection is the single point $x_b$, and $[T^3]\cdot[\Sigma^+_W]=\pm1$.
\item $\Sigma_q\cap\Sigma^+_W=\emptyset$: $\Sigma_q$ lies in the Morse polydisc, where $\tilde\chi_1$ and $\tilde\chi_2$ are within
$O(r)$ of $-\mu_P^{-1}<0$, while on $\Sigma^+_W$ they are positive.
\end{itemize}
Hence $\bar n=0$. For a point $p_0$ of the regular part and paths $\gamma_p$ from $p_0$ to $p$ in the connected set
$B_Y\setminus\Gamma^M$, the union of the fibre tori over $\gamma_p$, built region by region, is a four-chain with boundary
$T^3_p-T^3_{p_0}$. The chain $\check\Xi^W(z):=\check\Xi_0-\sum_p(\text{correction in }W^*_{V_p})-\sum_pn_p(\text{union of the
tori over }\gamma_p)$ has boundary $\sum_q\rho_q\Sigma_q-\bar nT^3_{p_0}=\sum_q\rho_q\Sigma_q$.
\end{proof}

The positive real locus thus replaces a local analysis of the fibre chains at the vertices of $\mathfrak T$: at the interior
vertices of a strict tropical $2$-cycle, at the branching vertices of the remainder, and at the points where two
summands that cross in $B_Y$ reach an identified curve.

\subsection{Assembly of the local lifts}\label{ssec:lift-assembly}

We now combine the local and regular chains and transport the result to the holomorphic hypersurface.
The preceding fibre-class argument removes the residual boundary, leaving exactly the prescribed sum of vanishing spheres.

\begin{remark}[where summands meet]\label{rem:summands-meet}
At a junction where a rotated exit strip reaches a reference side already carrying a strip, all adjacent coefficients
are multiples of $d^*_A$ and sum to zero. The pieces, collars and regular triangles therefore cancel
(Lemmas~\ref{lem:Mbd}, \ref{lem:collarB}(3) and~\ref{lem:Xichain}). This includes the three-way junction at the end of a
stretch. A germ and its own exit strip, the start of a rotation, and the half-edge of a tripod are interior to one
summand; their common boundary is treated within that summand.

Both passages at a node may occur. They give separate node pieces with the same sphere $\Sigma_q$ and total boundary
$(c_{ab}-c_{bc})\Sigma_q$ (Lemma~\ref{lem:nodesum}). Their halves meet along the half-axis $C^M(u',P)$ off
$\Gamma^M$, but remain separate in $\mathscr S$. The same convention applies when a crossing disc meets a sheet:
intersection of their images in $W^*$ does not identify their domains.

The remainder $y^\natural$ meets the germs and discs along $\partial U$, along replacement boundaries, and along the
seam $\{\varrho_{\mathsf P}=\varrho_0\}$ of a smoothed rotation. These are the identified curves of type (ii) in
$\mathscr S$, at distance at least $r_\natural$ from $\Gamma^M$; the seam coefficients are multiples of $d^*_A$.
Lemma~\ref{lem:Xichain} gives the cycle condition at their edges. At vertices, the remaining boundaries are closed
three-cycles in the fibre complexes $F_p$. Proposition~\ref{lem:F} removes their total fibre class by pairing with the
positive real locus. On the resolution side the proof instead uses the local fundamental groups
(Section~\ref{sec:resolution}).
\end{remark}

\subsubsection*{Order of choices}
The constants and choices are made in the following order.
\begin{enumerate}
\item The configuration: $\Delta$, $\mathcal T$ and $\check h$, $\mathcal F_C$ and $h_C$, the vectors $w_u$, and $k_0$ subject
to~\eqref{eq:lambdagap}; hence $\theta_0,\dots,\theta_3$, the gaps, the constants $\mathcal E_P$, the interior polygons
$\Omega_e$ and $C_\eta$; then the model positions, the model arc systems, $\theta_4$, $\Upsilon^M$, the charts $T_q$, $T_e$,
the constants $r_S$, $\varsigma_{\max}$ and $C_T$.
\item The model germs: the points $y^*_{A,0},y^0_A,\mathfrak x_A,y^1_A,y^*_{A,1}$ in $A\cap\Omega_e$ on every arc, then
$r_0$, at most $r_S/4$ and at most a quarter of the mutual distances of these points, of their distances from
$A\setminus\Omega_e$ and of their distances from the ends of $A$; the polygon $\mathsf P$, the unit germs with their orientations, the reference sides and the normal
discs; the balls $B^0_p$ of radius $r_B\le r_0$ and the constant $c_B$ of Section~\ref{ssec:lift-collar}; the orientations
of the linkage classes; the maps $\mathfrak h$, $\kappa_{\mathfrak h}$.
\item A compact set $\mathscr C_0\subset\mathcal Z$, $r\le r_1$ (Lemma~\ref{lem:NM}), $\eta_0\le r^2/20$,
$t<t_0(r,\eta_0,\mathscr C_0)$ as in Lemmas~\ref{lem:W}, \ref{lem:NM} and~\ref{lem:TR}, and $c$ as at the beginning of this
section.
\end{enumerate}
These do not depend on $z$. Then, for the given $z\in Z(\mathfrak K)$:
\begin{enumerate}
\setcounter{enumi}{3}
\item $\mathfrak K_1$ and $z^M=\mathfrak h_*z$; $\mathfrak K_2^0$, $\mathfrak K_2$, $\mathfrak K_2'$ and $U$; $w$, $v_A$ and $y$
by~\eqref{eq:nf1}; the local replacements with $\varepsilon_A$, $\varepsilon'_A$ and $\varrho_0$; $z^\natural$ and $r_\natural$
(Lemma~\ref{lem:nfY});
\item $\varsigma_0\le\min(\varsigma_{\max}/2,r_S/4,r_\natural/(4C_T),c_Br_B/8)$, so that $2\varsigma_0$ is small against the
radius of the balls $B^0_p$; the complex $\mathscr S$ and its triangulation
$\mathfrak T$, the regular sets $V_\kappa$ and $\epsilon$;
\item $t'\le\min(t,t(z))$, where $t(z)$ is a threshold for which Lemmas~\ref{lem:R1} and~\ref{lem:interface} hold for the sets
$V_\kappa$: the errors $O(1/L)$ of Lemma~\ref{lem:interface} must be small compared with the margins of the $V_\kappa$;
\item at $t'$: the pieces of depth $\varsigma_0$, the collar paths, the fibre complexes and the fillings $\mathscr Q_e$, the map
$F^\vee$, and the chain $\check\Xi^W(z)$ of Proposition~\ref{lem:F}.
\end{enumerate}
In particular $\varsigma_0$ is chosen after $U$ and $r_\natural$, and $t(z)$ after $\mathfrak T$; the threshold $t_0$ is that of
Lemmas~\ref{lem:W}--\ref{lem:TR} and does not depend on $z$.

\begin{proposition}[the lift, precise form]\label{prop:Bprecise}
Let $\Delta$ be admissible with a profile $\sigma$ satisfying~\ref{hyp:P}, and fix the configuration as in
Section~\ref{ssec:member}. Let $\mathscr C_0\subset\mathcal Z$ be compact, $r\le r_1$, $\eta_0\le r^2/20$,
$t<t_0(r,\eta_0,\mathscr C_0)$, and let $c$ be a real coefficient vector with the sign pattern, in the window of
Definition~\ref{def:window}. Then $Y_t$ is smooth; each node
parallelogram $P$ has a unique Morse point $q_P$ of the defining function in $\mathbb D_P$, with vanishing sphere $\delta_q$
oriented by Definition~\ref{def:orient}; and for every triangulation $\mathfrak K$ of $B'_\sigma$ in which $\Gamma$ is full
and every relative tropical $2$-cycle $z\in Z(\mathfrak K)$ there is an integral four-chain $\check\Xi(z)$ on $Y_t$,
constructed from $z$ and the choices listed above, with
\[
\partial\check\Xi(z)=\sum_q\rho_q(z)\,\delta_q .
\]
\end{proposition}

\begin{proof}
Smoothness of $Y_t$ and the Morse points are Lemmas~\ref{lem:W}(1) and~\ref{lem:NM}. Let $t'\le\min(t,t(z))$. The transfer,
the normal form and the choices of steps 4 and 5 do not depend on $t$. At $t'$, the pieces, the collar paths, the domain
and the extension give the chain $\check\Xi_0$, and Proposition~\ref{lem:F} gives $\check\Xi^W(z)$ on $W^*$ with boundary
$\sum_q\rho_q(z)\Sigma_q$. Transport it by the diffeomorphism $\Psi_1\colon W^*\to Y_{t'}$ of
Proposition~\ref{prop:isotopy}. By Lemmas~\ref{lem:NQ}(4) and~\ref{lem:carry}, $\Psi_1(\Sigma_q)$ is isotopic to
$\delta_q(t')$ in $Y_{t'}$, with the orientations of Definition~\ref{def:orient}; adding the traces of these isotopies gives
a chain on $Y_{t'}$ with boundary $\sum_q\rho_q(z)\delta_q(t')$. Finally the members $Y_{t''}$, $t''\in[t',t]$, with the same
$c$ form a smooth family of smooth hypersurfaces (Lemma~\ref{lem:W}(1) at $s=1$ for each $t''$), and the Morse points
$q_P(t'')$ vary smoothly with critical values bounded away from $0$ (Lemma~\ref{lem:NM}(3)); so the isotopy of the family
carries $\delta_q(t')$ to $\delta_q(t)$ up to isotopy, with orientation. The image of the chain is $\check\Xi(z)$.
\end{proof}

\begin{proof}[Proof of Theorem~\ref{thm:lift}]
The members $Y_t$ of Theorem~\ref{thm:main} are those of Proposition~\ref{prop:Bprecise}: with the constants
$\eta_1=r_1^2/20$ and $C=C_\eta$ of Section~\ref{ssec:member}, take $r=r_1$, so that $\eta_0\le\eta_1=r^2/20$; every real $c$
with $c_0>0>c_n$ is $c^{(0)}e^\xi$ for a unique $c^{(0)}\in\mathcal Z$ and a unique $\xi\in\operatorname{span}_\R\{r_P\}$
(Section~\ref{ssec:mirrorfamily}), with $\eta_P=\log\mu_P(c)$, so the conditions of Theorem~\ref{thm:main} on $c$ say that
$c$ is in the window. So Proposition~\ref{prop:Bprecise} gives $\check\Xi(z)$ with
$\partial\check\Xi(z)=\sum_q\rho_q(z)[\delta_q]$ for every $z\in Z(\mathfrak K)$. For a tropical $2$-cycle $(S,j,v)$ in the
sense of \cite[Def.~7.2]{CBM}, triangulated so that it is a relative tropical $2$-cycle (Section~\ref{ssec:cbm}), the chain
$\check\Xi(S)$ realises $\widetilde S$ as stated in Remark~\ref{rem:CBMlift}.
\end{proof}

\subsection{What the construction uses, and the strict case}\label{ssec:lift-cbm}

The proof uses from a relative tropical $2$-cycle only three pieces of data: its boundary network, which fixes the
multiplicities of the model germs; the coefficients $v_A$ of the transverse discs, which are arbitrary monodromy-invariant
covectors; and the remainder off a neighbourhood of the discriminant, which is an arbitrary chain with coefficients in the
local system. The strict tropical $2$-cycles of \cite[Def.~7.2]{CBM} impose conditions that the construction does not use,
as follows. The embeddedness of $S$ (clause (a)) and the condition that $S$ meets the discriminant only in $\partial S$ and
in finitely many crossings (clause (b)) are replaced by the normal form: the germs are embedded surfaces by construction,
and summands that cross in $B_Y$ are lifted separately over $\mathscr S$. Primitivity of the field (clause (c)) is
replaced by pieces and tori with integer multiplicities. The boundary vertices at the negative vertices of $B'_\sigma$
(clause (d)) become tripods with any multiplicity (Lemma~\ref{lem:unitgerms}), and the node passages (clause (e)) become
superpositions of the two unit halves, with both passages allowed (Lemma~\ref{lem:nodesum}). The trivalent interior edges
with a common rank-one sublattice (clause (g)) are replaced by fillings $\mathscr Q_e$ that exist for any valency. The condition
$|k_1|=1$ at the crossings (clause (i)) is replaced by monodromy invariance of the disc coefficient
(Lemma~\ref{lem:crossing}), and the interior vertices (clause (j)) are not used, on either route: Proposition~\ref{lem:F} disposes
of the fibre cycles there. The one condition that is used is the boundary condition (clause (h)): along the discriminant the
coefficient is a multiple of the vanishing covector $d^*_A$, so that the torus of a sheet along an arc is the exact torus
$\mathbb T_d$, which collapses to a circle over the arc and gives a piece without boundary there
(Lemma~\ref{lem:pieces}(1)). This is the relative condition $\mathcal D=\Z d$ of Definition~\ref{def:relcycle}, and the
monodromy invariance of the coefficients along the discriminant, which is part of the definition of $\iota_*\Lambda$, is
what the crossing pieces use (Lemma~\ref{lem:crossing}). Relative tropical
$2$-cycles are therefore the natural class for the lift: they are the chains whose local lifts exist, they form a group on
which the boundary $\sum_q\rho_q(z)[\delta_q]$ of the lift is additive, and their node coefficients already exhaust $K$ (Theorem~\ref{thm:real}).

\begin{remark}[the strict case]\label{rem:CBMlift}
Let $(S,j,v)$ be a tropical $2$-cycle in the sense of \cite[Def.~7.2]{CBM} and $z=[S]$ its triangulation.

\emph{Away from the discriminant.} Outside $U$ the normal form is the transfer of $S$ itself: $y$ agrees there with
$\operatorname{sd}z^M$. Over that part the domain is the bundle of the cells $\kappa\times\mathbb T_\kappa$ with $m_\kappa=1$
and $\mathbb T_\kappa=\ker\kappa_{\mathfrak h}(v)\otimes U(1)$, the fibre of the object $\widetilde S$ that Casta\~no-Bernard and Matessi
attach to $S$ \cite[\S7, (46)]{CBM}, written in the affine structure of $B_Y$ through the flat pairing of
Proposition~\ref{prop:L3} (Remark~\ref{rem:CBMfibre}). The map $F^\vee$ is built cell by cell in the regions
$W^*_{V_\kappa}$ (Lemma~\ref{lem:extB}), so over this part the chain is the bundle of the two-tori $\ker\kappa_{\mathfrak h}(v)\otimes U(1)$
up to homotopy in each region.

\emph{Near the discriminant.} Since $\partial S$ is an embedded curve passing each vertex of the discriminant at most once,
the multiplicities of the model germs are $0$ or $\pm1$. The chain agrees with the local models of \cite{CBM} under the
following two conditions:
\begin{enumerate}
\item[(C1)] every arc of $\Gamma^M$ carries at most one crossing of the transfer of $S$;
\item[(C2)] within $U$ the transfer of $S$ is the sum of the model germs (M1)--(M4) with these multiplicities and of normal
discs at its crossings: at a model node it is the standard half on the side $\alpha_q\ge0$, at a model vertex the tripod on
the side of $u'$, and along each arc of $\partial S$ it lies on the model sides with the windings of the model germs.
\end{enumerate}
Under (C1) and (C2), the part of $\operatorname{sd}z^M$ on $U$ minus $\mathrm{Std}^M(b^M)\cap U$ is a sum of normal discs at the crossings, homologous relative to
$\partial U$, by moving discs along the arcs, to $\sum_Av_A\bar D(E_A)$ with $v_A$ the crossing coefficient of $S$ on $A$
(and $v_A=0$ on arcs without crossing); so the normal form may be taken with these $v_A$, which are primitive with
$|k_1|=1$ \cite[Def.~7.2(i)]{CBM}, and all pieces are the local models of \cite{CBM} listed below. Without (C1) or (C2) the
transverse disc on an arc collects the crossings on that arc and the differences of sides and windings, and its piece is
that of Lemma~\ref{lem:crossing} for the collected coefficient, which need not be a local model of \cite{CBM}: two crossings
with $k_g=1$ on one arc give $k_g=2$, a two-sheeted cover with a fibre of type $I_2$; a sheet of $S$ on the side $\alpha_q\le0$ at
a node gives a coefficient in $\Z d^*_A$ and the piece $\mathbb T_d\cdot\mathfrak s^\vee(\mathfrak D_{\varsigma_0})$.

Condition (C2) refers to conventions fixed before $z$ (step~2 of the order of choices): the basis
$(\mathrm{Turn}(u_1),\mathrm{Turn}(u_2))$ at each turning point, the side $\alpha_q\ge0$ at the model nodes and the side of
$u'$ at the model vertices, and the reference side $\alpha_e\ge0$ of each arc. With these fixed conventions (C2) fails for
many strict cycles: for instance when $\partial S$ turns through the sector $C(u_3)$, which has multiplicities $(-1,1)$ in
the basis $(\mathrm{Turn}(u_1),\mathrm{Turn}(u_2))$, or when a sheet of $S$ lies on the side $\alpha_q\le0$ or on the
non-reference side of an arc. Any two of $\mathrm{Turn}(u_1),\mathrm{Turn}(u_2),\mathrm{Turn}(u_3)$ form a basis, and the
other choices of side and of reference side change the normal form only by transverse discs, so the construction and
Theorem~\ref{thm:lift} hold for every choice of these conventions, and they may be chosen depending on $S$. If the
conventions are chosen as those used by $S$ (at each turning point a basis containing the turning
germ of $\partial S$, at each model node and model vertex the side of $S$, on each arc the side of $S$ as reference
side), and if (C1) holds and within $U$ the transfer of $S$ is the sum of the germs so chosen and of normal discs at its
crossings, then all pieces are the local models of \cite{CBM}. This comparison assumes the stated position of $S$.
\begin{center}\small
\begin{tabular}{p{5.2cm}p{4.3cm}p{4.6cm}}
part of $\check\Xi(S)$ & constructed in & counterpart in \cite{CBM}\\\hline
over the sheets: torus bundle & Section~\ref{ssec:lift-regular} & \cite[\S7, (46)]{CBM}\\
over interior edges & Section~\ref{ssec:lift-regular} & \cite[\S7.7, Step~2]{CBM}\\
node piece & Lemmas~\ref{lem:NQ}, \ref{lem:nodesum} & \cite[Lemma~5.14]{CBM}\\
pieces at the model vertices & Lemma~\ref{lem:pieces}(2) & \cite[\S7.5; \S7.7, Steps~2--3]{CBM}\\
pieces at turning points & Lemma~\ref{lem:pieces}(3) & \cite[Lemma~7.6]{CBM}\\
pieces at crossings & Lemma~\ref{lem:crossing}(2) & \cite[\S7.7, Step~6]{CBM}\\
interior vertices & Proposition~\ref{lem:F} & replaces \cite[\S7.7, Steps~4--5]{CBM}
\end{tabular}
\end{center}
We do not use that $\widetilde S$ is a manifold at the interior vertices of $S$ \cite[\S7.7, Step~5]{CBM}.
\end{remark}
